% concatenated from main.tex
\documentclass{article}


% if you need to pass options to natbib, use, e.g.:
%     \PassOptionsToPackage{numbers, compress}{natbib}
% before loading neurips_2025

% ready for submission
\usepackage[final]{neurips_2025}
% \usepackage{neurips_2025}
\usepackage[center]{caption}
\usepackage{algorithm}
\usepackage{algpseudocode}
\usepackage{geometry}
\usepackage{mathtools}
\usepackage{float}
\usepackage{caption} % dans le préambule
\usepackage{titletoc}

\geometry{margin=1in}
\usepackage{amsthm}
\usepackage{float}
\usepackage{caption}
\usepackage{graphicx}
\newtheorem{theorem}{Theorem}[section]
\newtheorem{proposition}[theorem]{Proposition}
\newtheorem{lemma}[theorem]{Lemma}
\newtheorem{corollary}[theorem]{Corollary}
\newtheorem{definition}[theorem]{Definition}
\newtheorem{remark}[theorem]{Remark}
\setlength{\abovedisplayskip}{3pt}
\setlength{\belowdisplayskip}{3pt}
\setlength{\abovedisplayshortskip}{3pt}
\setlength{\belowdisplayshortskip}{3pt}
\usepackage{algorithm}
\usepackage{algpseudocode}
% to compile a preprint version, e.g., for submission to arXiv, add add the
% [preprint] option:
%     \usepackage[preprint]{neurips_2024}


% to compile a camera-ready version, add the [final] option, e.g.:



% to avoid loading the natbib package, add option nonatbib:
%    \usepackage[nonatbib]{neurips_2024}

\usepackage{wrapfig} % Package for wrapping text around figures
\usepackage[utf8]{inputenc} % allow utf-8 input
\usepackage[T1]{fontenc}    % use 8-bit T1 fonts
\usepackage{hyperref}       % hyperlinks
\usepackage{url}            % simple URL typesetting
\usepackage{booktabs}       % professional-quality tables
\usepackage{amsfonts}       % blackboard math symbols
\usepackage{nicefrac}       % compact symbols for 1/2, etc.
\usepackage{microtype}      % microtypography
\usepackage{xcolor}         % colors

\usepackage{amsmath, amssymb}
% to compile a preprint version, e.g., for submission to arXiv, add add the
% [preprint] option:
%     \usepackage[preprint]{neurips_2025}


% to compile a camera-ready version, add the [final] option, e.g.:
%     \usepackage[final]{neurips_2025}


% to avoid loading the natbib package, add option nonatbib:
%    \usepackage[nonatbib]{neurips_2025}


\usepackage[utf8]{inputenc} % allow utf-8 input
\usepackage[T1]{fontenc}    % use 8-bit T1 fonts
\usepackage{hyperref}       % hyperlinks
\usepackage{url}            % simple URL typesetting
\usepackage{booktabs}       % professional-quality tables
\usepackage{amsfonts}       % blackboard math symbols
\usepackage{nicefrac}       % compact symbols for 1/2, etc.
\usepackage{microtype}      % microtypography
\usepackage{xcolor}         % colors

\bibpunct{[}{]}{,}{n}{}{,}
\bibliographystyle{plainnat}


\title{How Patterns Dictate Learnability in Sequential Data}


% The \author macro works with any number of authors. There are two commands
% used to separate the names and addresses of multiple authors: \And and \AND.
%
% Using \And between authors leaves it to LaTeX to determine where to break the
% lines. Using \AND forces a line break at that point. So, if LaTeX puts 3 of 4
% authors names on the first line, and the last on the second line, try using
% \AND instead of \And before the third author name.


\author{
  Mario Morawski\thanks{Equal Contribution} \\
  % \texttt{mario.morawski@free.fr} \\
  \And
  Anaïs Després\footnotemark[1] \\
  % \texttt{andespres@gmail.com} \\
  \AND
  Rémi Rehm \\
}


\begin{document}


\maketitle


\begin{abstract}
Sequential data—ranging from financial time series to natural language—has driven the growing adoption of autoregressive models. However, these algorithms rely on the presence of underlying patterns in the data, and their identification often depends heavily on human expertise. Misinterpreting these patterns can lead to model misspecification, resulting in increased generalization error and degraded performance. The recently proposed \texttt{evolving pattern (EvoRate)} metric addresses this by using the mutual information between the next data point and its past to guide regression order estimation and feature selection. Building on this idea, we introduce a general framework based on predictive information—the mutual information between the past and the future, $\mathbf{I}(X_{\text{past}}; X_{\text{future}})$. This quantity naturally defines an information-theoretic learning curve, which quantifies the amount of predictive information available as the observation window grows. Using this formalism, we show that the presence or absence of temporal patterns fundamentally constrains the learnability of sequential models: even an optimal predictor cannot outperform the intrinsic information limit imposed by the data. We validate our framework through experiments on synthetic data, demonstrating its ability to assess model adequacy, quantify the inherent complexity of a dataset, and reveal interpretable structure in sequential data.

\end{abstract}

\section{Introduction}
From time series in finance and healthcare \cite{ashok2023tactis, rasul2020multivariate, lim2021time} to text streams in natural language processing \cite{naveed2023comprehensive}, much of the real-world data ingested by machine learning systems is inherently sequential. Autoregressive models are commonly trained to predict such data, with their performance relying heavily on capturing evolving patterns. However, identifying these patterns still depends largely on human expertise \cite{zeng2025towards}. Misinterpretation can result in models that overfit to training data, leading to poor calibration, generalization errors, and heightened vulnerability to adversarial examples \cite{fischer2020conditional}; for instance, classifiers can memorize random label assignments in the training set \cite{zhang2016understanding}.
A substantial body of work is devoted to bounding generalization error, primarily through Bayesian theory \cite{alquier2024user}. For sequential data, these bounds have been formulated using Rademacher complexity \cite{mcdonald2011rademacher}, leveraging concentration inequalities and Bayesian-inspired methods.
We identify a critical open question in the study of sequential data: \textbf{(i) What is the minimal achievable risk for a predictor attempting to model sequential data?} This question differs subtly but importantly from bounding the gap between empirical and true risk. Empirical evidence, such as the plateau in model performance on the \emph{Exchange} dataset from \texttt{GluonTs} \cite{alexandrov2020gluonts} despite recent innovations \cite{rasul2021autoregressive}, suggests that data limitations—rather than model inadequacy—may be the binding constraint.
This leads to a critical second question:
\textbf{(ii) Can we distinguish whether poor performance stems from the model’s limitations or from the inherent unpredictability of the data?}
This distinction hinges on the nature and strength of temporal patterns: fewer patterns imply that even an optimal predictor will perform poorly, while richer structure offers more room for effective forecasting. Thus, good prediction requires both $(1)$ that the data contain exploitable patterns, and $(2)$ that the model can identify and leverage them.
To illustrate this, consider the following stylized scenarios in which a colleague attempts to predict what meal a researcher will bring to lunch each day:
\begin{itemize}
    \item Meals are sampled uniformly at random from a cookbook.
    \item Meals are chosen weekly based on preferences, weather, and prior meals.
    \item Meals follow a fixed rotation established years ago.
\end{itemize}
In the first case, even the optimal predictor cannot outperform random guessing due to the absence of structure. In contrast, the second and third cases involve latent patterns—periodicity, preference dependencies, contextual triggers—that a well-designed model could exploit. Crucially, identifying whether such structure exists—and how much of it is learnable—is a prerequisite for developing effective models.
Concurrently, recent work has pointed out the lack of effective metrics for quantifying evolving patterns in sequential data, proposing the \texttt{evolving rate (EvoRate)} metric as a preliminary solution \cite{zeng2025towards}. Building on this idea, we generalize the notion of \texttt{EvoRate} through the concept of \emph{predictive information}—the mutual information between the past and the future \cite{bialek1999predictive, crutchfield1997statistical}. This leads to the formulation of the \textbf{universal learning curve} \cite{bialek1999predictive}, which measures the rate at which predictive information increases with observation length. This curve serves as a fundamental tool to quantify the minimal achievable risk. Our framework thus provides an information-theoretic perspective that addresses both questions \textbf{(i)} and \textbf{(ii)}, while also offering insights into the nature and strength of the patterns embedded in the data.
Our main contributions are:
\begin{itemize}
\item \textbf{Establishing} a theoretical link between the presence of temporal patterns in sequential data and the minimal achievable prediction risk.
\item \textbf{Deriving} information-theoretic bounds on this minimal risk, expressed in terms of structural properties of the data such as predictive information and pattern complexity.
\item \textbf{Proposing} a practical estimator of the intrinsic risk limit, enabling quantitative comparison with model performance to assess whether learning is limited by data or by the model.
\item \textbf{Validating} the framework on synthetic data, showing how it supports model selection and reveals when further improvements are constrained by the data itself. The codes are available on GitHub: \url{https://github.com/EkMeasurable/Learnability_Ipred}
\end{itemize}


\section{Related work}

\paragraph{Patterns Estimation for Sequential Data}
Early work on measuring predictability in time series includes \texttt{ForeCA} \cite{goerg2013forecastable}, which introduced the concept of forecastability, based on the entropy of the spectral density. This metric quantifies how predictable a time series is from a frequency-domain perspective. More recently, \texttt{EvoRate} \cite{zeng2025towards} proposed a mutual information-based approach to capture evolving patterns in sequential data, and demonstrated its use for guiding regression order and feature selection.
While promising, both \texttt{ForeCA} and \texttt{EvoRate} focus primarily on intrinsic signal properties and lack a principled connection to model performance. In particular, \texttt{EvoRate} emphasizes temporal changes in its value rather than its absolute magnitude, which remains under-exploited. Moreover, these methods do not assess whether a predictive model has effectively captured the patterns present in the data.
In contrast, we introduce an information-theoretic estimator of the minimal achievable risk on a dataset, which serves two key purposes: (i) indicating the presence and strength of temporal patterns—lower risk suggests stronger patterns—and (ii) enabling direct comparison with a model's empirical risk. 
% A small gap between the two implies that the model has effectively exploited the available structure; a large gap suggests potential for improvement. This dual interpretability bridges the gap between pattern quantification and model evaluation.

\paragraph{Learning under General Stochastic Processes.}  
Recent advances in statistical learning theory have increasingly focused on relaxing the classical assumption of independent and identically distributed (i.i.d.) data. In particular, extending learning frameworks to handle stochastic processes with temporal dependencies or evolving distributions has become a central research direction. The \emph{Prospective Learning} framework~\cite{desilva2025prospectivelearninglearningdynamic} formalizes this setting by requiring a learner to produce a sequence of hypotheses that achieve low risk on future observations, given the data observed up to the present—an approach particularly relevant in non-stationary or dynamic environments where the optimal predictor may change over time. Complementary efforts have characterized learnability under general stochastic processes~\cite{dawid2022learnabilitygeneralstochasticprocesses, hanneke2020learninglearningpossibleuniversal}, providing general conditions for when consistent learning is possible, often framed in terms of regret minimization or universal consistency. These perspectives share the common goal of determining whether, and under what conditions, learning remains feasible in complex non-i.i.d. settings. Our work complements these approaches by introducing an intrinsic, information-theoretic notion of learnability. Rather than relying solely on external performance measures such as risk or regret, we quantify learnability through the predictive structure of the data itself—captured by predictive information and the universal learning curve. This allows us to bridge theoretical learnability guarantees with the structural properties of the data-generating process, offering a complementary perspective on when and why generalization is possible in sequential environments.


\paragraph{Mutual Information Estimation (MI)}
Traditional methods, such as the k-nearest neighbor estimator \cite{kraskov2004}, perform well in low-dimensional settings but struggle when applied to long or high-dimensional time series. MINE \cite{belghazi2018} addresses these challenges by leveraging the Donsker–Varadhan representation and deep neural networks to learn flexible estimators of mutual information. However, MINE exhibits high variance and instability, particularly in sequential data settings. Alternative approaches, such as InfoNCE \cite{oord2018}, originally developed for contrastive predictive coding (CPC), are better suited to time series tasks; they learn representations that maximize mutual information between past and future segments. Further advancements include CLUB \cite{cheng2020}, which provides a tractable upper bound on mutual information, and SMILE \cite{song2019}, which improves estimator stability through Jensen–Shannon–based objectives.

\paragraph{Predictive Information and Learning Curve.}
Originally introduced by \cite{bialek1999predictive, bialek2001predictability}, predictive information can be viewed as a generalization of \texttt{EvoRate}. The core idea is to measure the mutual information between a context, $X_{\text{past}}$, and a target, $X_{\text{future}}$. It has been applied in machine learning to learn meaningful data representations \cite{lee2020predictive}, and several methods for estimating it have been proposed \cite{fischer2020conditional}.
Building on this line of work, \cite{bialek1999predictive} and \cite{crutchfield2003regularities} also introduced the notion of the \emph{universal learning curve}, which describes how predictive information grows with the length of the observed context. This curve captures the fundamental limits of learnability by quantifying how much additional information about the future can be extracted from longer histories.
In this paper, we extend these theoretical ideas by establishing a novel link between the universal learning curve and the \emph{minimal achievable risk}. Specifically, we show that the learning curve can be used to derive an empirical estimator of the optimal prediction performance attainable by any model. This connection provides an operational interpretation of predictive information as a diagnostic tool to distinguish between model limitations and the intrinsic unpredictability of the data.

\paragraph{Lowest Possible Error Rate}
In classification tasks, the Bayes error rate represents the lowest achievable error rate for any classifier \cite{tumer1996estimating, chen2023evaluating, tumer2003bayes}. This concept was developed alongside the PAC-Bayes framework, which provides upper bounds on the risk \cite{rivasplata2019pac, langford2001bounds}. In the context of sequential data, prior work has focused primarily on bounding the gap between empirical and true risk, notably through Rademacher complexity \cite{mohri2008rademacher, mcdonald2011rademacher, kuznetsov2015learning}. However, to the best of our knowledge, none of these approaches explicitly connects the minimal achievable risk to the presence of patterns in sequential data.

\section{Preliminary}
\subsection{Notations and hypothesis} 

We consider a stochastic process $\mathbf{X}_{t}^{T} = \{X_u\}_{u=t}^{T} \in \mathcal{X}^{T-t+1}$, indexed over $[t, T]$, with $t \in [0, T]$. We assume:
\[
\mathbf{(H_0)}: \quad \text{The process is stationary and } H(X_t) < \infty,
\]
where $H(\cdot)$ denotes entropy. We do not require $\mathbb{E}(X_t^2) < \infty$, allowing for heavy-tailed distributions, and typically set $\mathcal{X} = \mathbb{R}^d$.

Given a predictor $g: \mathcal{X}^{k} \to \mathcal{X}$ and a loss function $\ell: \mathcal{X} \times \mathcal{X} \to \mathbb{R}$, we define the \textbf{forecasting risk} of order $k$ as:
\begin{equation}
\mathcal{R}^{(k)}(g) = \mathbb{E}_{\mathbf{X}_{t-k+1}^{t+1}} \left[ \ell(X_{t+1}, g(\mathbf{X}_{t-k+1}^{t})) \,\middle|\, \mathbf{X}_{t-k+1}^{t} \right].
\end{equation}

\subsection{Entropy and Entropy Rate}

Understanding the nature of sequential data is closely linked to studying the memory properties of the underlying process. For a stationary process $\mathbf{X}_{t}^{T}$, a fundamental quantity for measuring the information content is its entropy \cite{shannon1948mathematical}, defined as
\begin{equation}
    H(\mathbf{X}_{t}^{T}) = - \int_{S} p(\mathbf{X}_{t}^{T}) \ln p(\mathbf{X}_{t}^{T}) \, d\mathbf{X},
\end{equation}
where $S$ denotes the support of the distribution. Under Assumption $\mathbf{(H_0)}$, the joint density $p(x_t, \dots, x_T)$ is invariant under time translation, so the entropy depends only on the length of the block $[t, T]$. The entropy is concave and subadditive with respect to the block size, and satisfies $H(0) = 0$. Let $l(k) = H(X_{t} \mid \mathbf{X}_{t-k+1}^{t-1}) = -\mathbb{E}_{P(X_{t}, \mathbf{X}_{t-k+1}^{t-1})} \ln P(X_{t} \mid \mathbf{X}_{t-k+1}^{t-1})$ be the \textbf{entropy rate} of order $k$ \cite{cover1999elements}. It measures how well $X_{t}$ can be predicted by observing $k$ past values. Naturally, the more past observations are available, the lower the uncertainty, so the mapping $k \mapsto l(k)$ is non-increasing and non-negative. 
Under Assumption $\mathbf{(H_0)}$, $l(k)$ converges to a constant $l_0$ as $k \to \infty$, and the \textbf{fundamental theorem of entropy} \cite{shannon1948mathematical} states that
\begin{equation}
\frac{1}{k} H(\mathbf{X}_{t-k+1}^{t}) \underset{k\to\infty}{\longrightarrow} l_0,
\end{equation}
that is $H(k) \sim k l_0$ when $k$ is large.

\subsection{$\mathbf{I}_{\text{pred}}$ as a generalization of \texttt{EvoRate}}
In~\citet{zeng2025towards}, \texttt{EvoRate} is proposed as a measure for capturing predictive patterns in sequential data. It is defined as the mutual information between the past $k$ observations, $\mathbf{X}_{t-k+1}^{t}$, and the next observation, $X_{t+1}$, that we note $\mathbf{I}(\mathbf{X}_{t-k+1}^{t}; X_{t+1})$.
We naturally extend \texttt{EvoRate} by introducing a fundamental quantity,
\begin{equation}\label{eq:evo_pred_def}
\mathbf{I}_{\text{pred}}(k,k^\prime) = \mathbf{I}(\mathbf{X}_{t-k+1}^{t}; \mathbf{X}_{t+1}^{t+k^\prime}) = \int p(\mathbf{X}_{t-k+1}^{t+k^\prime}) \ln \frac{p(\mathbf{X}_{t-k+1}^{t+k^\prime})}{p(\mathbf{X}_{t-k+1}^{t})p(\mathbf{X}_{t+1}^{t+k^\prime})} d\mathbf{X}.
\end{equation}
Predictive information, denoted $\mathbf{I}_{\text{pred}}$, quantifies how much information from the past can be used to predict the future. Under assumption $\mathbf{(H_{0})}$, it converges in the limit as both context lengths grow:
\begin{equation}
\lim_{k, k^\prime \to \infty} \mathbf{I}_{\text{pred}}(k, k^\prime) = \mathbf{I}_{\text{pred}}(X_{\text{past}}, X_{\text{future}}),
\end{equation}

a quantity also known in the literature as \emph{excess entropy}~\cite{crutchfield1983symbolic, crutchfield1997statistical, feldman1998discovering} or the \emph{effective measure of complexity}~\cite{grassberger1986toward}. Our focus remains on finite $k$ and $k^\prime$ values to ensure practical relevance.
Under $\mathbf{(H_{0})}$ and suitable regularity conditions, predictive information corresponds to the sub-extensive component of the entropy.

The asymptotic behavior of predictive information provides a principled lens through which to assess the complexity and predictability of stochastic processes~\cite{bialek1999predictive, crutchfield2003regularities}. In machine learning, recent works have leveraged this quantity to improve learned representations~\cite{lee2020predictive, fischer2020conditional}, taking advantage of its ability to capture shared temporal structure across sequences.

\section{Theoretical analysis of $\mathbf{I}_{\text{pred}}$}

Consider the set $\mathcal{F}_k = \left\{ \left( \mathbf{X}_{0}^{t}, X_{t+1} \right) \mapsto -\ln Q\left( X_{t+1} \mid \mathbf{X}_{t-k+1}^{t} \right) \mid Q \in \mathcal{H}_k \right\}$, which defines a class of real-valued functions on $\mathbf{X}$. This class encompasses all possible loss functions that can be derived from making predictions using a model $Q \in \mathcal{H}_k = \left\{ Q_{\theta} \mid \theta \in \Theta \right\}$, where $\mathcal{H}_k$ represents a family of models of order $k$ parameterized by $\theta$—the vector of model parameters that governs the behavior and structure of the prediction model.

\subsection{From a learning perspective}

The $k^{\text{th}}$-order forecasting risk $\mathcal{R}^{k}(Q)$, defined as 
\begin{equation}
\mathcal{L}_{\mathrm{mle}}^{k} = -\mathbb{E}_{P(X_{t+1}, \mathbf{X}_{t-k+1}^t)} \ln Q(X_{t+1} \mid \mathbf{X}_{t-k+1}^t),
\end{equation}
has been widely used to assess autoregressive model performance~\cite{rangapuram2018deep, zeng2025towards}. In a similar spirit, we draw a connection between the predictive information $\mathbf{I}_{\text{pred}}$ and a central concept in information theory: the \emph{universal learning curve} $\Lambda(k) = \ell(k) - \ell_0$~\cite{bialek1999predictive, crutchfield2003regularities}, also known as the \emph{entropy gain}. This quantity measures the reduction in uncertainty about the future obtained by conditioning on $k$ past observations, and thus captures the presence of temporal patterns in sequential data.
A key theoretical connection between the predictive information and the learning curve is established by the following result from~\citet{bialek1999predictive}:
\begin{proposition}[\textit{Bialek and Tishby (1999)} \cite{bialek1999predictive}]\label{prop:Approximation_of_the_Universal_Learning_Curve}
Under hypothesis $\mathbf{(H_0)}$, we have:
\[
\mathbf{I}_{\text{pred}}(k+1,k') - \mathbf{I}_{\text{pred}}(k,k') \longrightarrow \Lambda(k) \quad \text{as} \quad k' \to \infty.
\]
\textit{Proof in Appendix~\ref{prop:learning_curve_approximation}.}
\end{proposition}
The notion of a learning curve is central to our analysis. In \texttt{EvoRate}~\cite{zeng2025towards}, the authors primarily focused on how the metric evolves with the size of the past window, as its absolute values were difficult to interpret. In contrast, the universal learning curve offers a more principled alternative: it can be seen as a discrete derivative of the predictive information \cite{bialek1999predictive}, capturing the marginal contribution of each additional past observation. This perspective not only enhances interpretability but also plays a key role in our theoretical developments and empirical evaluations.

\subsection{Interpreting the Asymptotic Behavior of $\Lambda(k)$}

The asymptotic behavior of $\Lambda(k)$ offers insight into the nature of the temporal patterns present in the data. It helps characterize the structure of the underlying distribution and supports the theoretical link between $\Lambda(k)$ and the presence of predictive patterns. We now present examples that illustrate the relevance of this perspective.

\paragraph{The special case of Markov processes.}

When \( X \) is generated by a Markov process of order \( m \), dependencies are limited to the past \( m \) observations. This structure is naturally reflected in the predictive information \( \mathbf{I}_{\text{pred}} \). Previous approaches have addressed Markov order estimation using conditional mutual information~\cite{papapetrou2013markov}, with \texttt{EvoRate} showing promising empirical results~\cite{zeng2025towards}, though without formal guarantees. Classical criteria like AIC~\cite{katz1981some} and BIC~\cite{csiszar2000consistency} also remain standard tools.

Building on~\citet{crutchfield2003regularities}, we derive a closed-form expression for \( \mathbf{I}_{\text{pred}} \) under the Markov assumption and show that the universal learning curve \( \Lambda(k) \) vanishes for all \( k \geq m \), providing a theoretical link between pattern complexity and memory length.

\begin{proposition}[Predictive information in Markov processes]\label{prop:markov_combined}
Let $\mathbf{X}_{t}^{T}$ be a Markov process of order $m$. If $k' \geq k \geq m$, then:
\begin{align*}
  \text{(i)}\quad & \mathbf{I}_{\text{pred}}(k,k')
  = \mathbb{E}_{\mathbf{X}_{t-m+1}^{t+m}}
  \left[\ln\frac{P\bigl(X_{t+1}^{t+m}\mid\mathbf{X}_{t-m+1}^t\bigr)}
                {P\bigl(X_{t+1}^{t+m}\bigr)}\right], \\[6pt]
  \text{(ii)}\quad & \forall\,k\ge m,\quad \Lambda(k) = 0.
\end{align*}
In particular, for first-order Markov processes, we recover that $\mathbf{I}_{\text{pred}}(k, k') = \texttt{EvoRate}(1)$ for all $k \ge 1$, allowing $\Lambda(k)$ to identify the true Markov order.\\
\textit{Proof in Appendices~\ref{prop:value_of_evoPred_for_markovian_process} and~\ref{proof:markovian_process_universal_learning_curve}.}
\end{proposition}

\paragraph{Parametric Processes.}

Let the process $\mathbf{X}_{t-k+1}^{t}$ be generated by a parametric family $Q_{\mathbf{X}_{t-k+1}^{t}}(\theta)$, with unknown parameter $\bar{\theta}$ drawn from prior $\mathcal{P}(\theta)$ and $\dim\Theta = p$. Classical results in the i.i.d.\ case \cite{bialek1999predictive, clarke1990information} show that predictive information grows logarithmically with $k$. We extend this to dependent settings under mild regularity conditions.

\begin{theorem}[Predictive information in parametric models]
\label{thr:evoPredLimitParametrizedProcess}
Assuming standard hypotheses on stationarity, weak dependence, and regularity of the parametric family, (\ref{H:ergodic}, ~\ref{H:ident}, and ~\ref{H:entropy}), let the past and future windows grow with $k\!\to\!\infty$ and
$k'\!\ge\!k$ (the ratio $k'/k$ may vary).
Then
\begin{equation*}\label{eq:evoPredLimitParametrizedProcess}
   \mathbf{I}_{\text{pred}}(k,k') \underset{k\to\infty}{=} \frac{p}{2}\ln(k) + \frac{1}{2} \ln\det(F) + \mathcal{O}(1),
\end{equation*}
where $F$ is the Fisher information matrix associated with the divergence $D_{KL}(Q_{\mathbf{X}_{t-k+1}^{t}}(\theta) \,\|\, Q_{\mathbf{X}_{t-k+1}^{t}}(\bar{\theta}))$.\\
\textit{Proof in Appendix~\ref{thm:PI}.}
\end{theorem}

This result provides an interpretable asymptotic expansion: the first term captures the uncertainty due to parameter estimation, while the second reflects model confidence via the concentration of likelihood around $\bar{\theta}$. In contrast to \texttt{EvoRate}, whose asymptotic behavior remains unclear, $\mathbf{I}_{\text{pred}}$ exhibits a principled structure grounded in classical learning theory.

\begin{corollary}[Universal Learning Curve Decay]
\label{cor:universal_learning_curve_parametric}
Under the same assumptions, the universal learning curve satisfies:
\[
\Lambda(k) \sim \frac{p}{2k}.
\]
In particular, the dimensionality of the parameter space can be estimated as $\dim \Theta \approx 2k\,\Lambda(k)$.\\
\textit{Proof in Appendix~\ref{thm:PI}.}
\end{corollary}

This relation reveals a precise connection between the decay of the universal learning curve $\Lambda(k)$ and the intrinsic dimensionality $p$ of the parameter space $\Theta$. The behavior is consistent with classical Bayesian learning theory in the i.i.d.\ case, now extended to structured and dependent sequences. This correspondence indicates that predictive information faithfully reflects the effective degrees of freedom of the generative process.

\paragraph{Beyond Finite Parametric Models.}

When the data-generating process depends on an infinite or unbounded set of latent parameters~\cite{bialek1999predictive, barron1991minimum}, predictive information may vanish or grow sub-logarithmically, indicating insufficient structure for reliable forecasting or a mismatch between model and data complexity. In such cases, structural assumptions or inductive biases may be necessary to enable generalization and avoid overfitting.

\subsection{Link Between Learning Curve and Minimal Achievable Risk}

We now show how the universal learning curve $\Lambda(k)$ provides an upper bound on how close a model of order $k$ can get to the optimal regression risk.
\begin{proposition}
For any $k \in \mathbb{N}$ and any $Q \in \mathcal{H}_k$,
\begin{equation*}
   \mathcal R^{\infty}(Q^{*})
   \;\le\;
   \mathcal R^{k}(Q) - \Lambda(k).
\end{equation*}
where $\mathcal{R}^{\infty}(Q^*)$ is the minimal risk achievable by the optimal predictor $Q^* = P(X_{t+1} \mid \mathbf{X}_{\text{past}})$.\\
\textit{Proof in Appendix~\ref{prop:B1}.}
\end{proposition}

In practice, the optimal achievable risk is given by $\mathcal{R}^{\infty}(Q^{*}) = \lim_{k \to \infty} H(X_{t+1} \mid \mathbf{X}_{t-k+1}^{t})$. This notion parallels the Bayesian risk in classification tasks~\cite{snapp1995estimating}, where the focus lies on the gap between the current and optimal losses, particularly in online learning contexts~\cite{hartline2015no, hsu2019empirical, lu2022no}.

Here, we aim to estimate $\mathcal{R}^{\infty}(Q^{*})$ directly. The previous proposition shows that no model with finite memory $k$ can attain the optimal risk unless $\Lambda(k)$ is negligible. The term $\Lambda(k)$ quantifies the irreducible uncertainty due to limited context and serves as a data-dependent upper bound on the possible reduction in risk achievable by optimizing over $\mathcal{H}_k$.

Since the true risk $\mathcal{R}^{k}(Q)$ is not directly observable, we use its empirical estimate,
\begin{equation}
\hat{\mathcal{R}}^{k}_{n}(Q) = \frac{1}{n} \sum_{i=1}^n -\ln Q\left(X_{i,t+1} \mid \mathbf{X}_{i,t-k+1}^{t}\right).
\end{equation}
To control the deviation between the true and estimated risks, we leverage Rademacher complexity for stationary sequences~\cite{mcdonald2011rademacher}, together with standard concentration inequalities.


\begin{proposition}
Assume $(X_t)_{t\in\mathbb{Z}}$ is a stationary process satisfying the conditions of \cite{mcdonald2011rademacher}. Then, for any $\delta\in(0,1)$, with probability at least $1-\delta$ over an $n$‑sample drawn from the process, every $Q\in\mathcal{H}$ satisfies
\begin{equation*}
    \mathcal{R}^{\infty}(Q^{*})  \leq \hat{\mathcal{R}}^{k}(Q)-\Lambda(k)+2 \widehat{\Re}_n\left(\mathcal{F}_k\right)+3 \frac{\ln (1 / \delta)}{n}
\end{equation*}
where 
\[
    \widehat{\Re}_n(\mathcal{F}_k)
    \;=\;
    \mathbb{E}_{\sigma}\!\Biggl[\sup_{f\in\mathcal{F}_k}
    \frac{1}{n}\sum_{i=1}^n \sigma_i\,f\bigl(\mathbf{X}_{i,0}^t,X_{i,t+1}\bigr)\Biggr]
\]
is the empirical Rademacher complexity of the class $\mathcal{F}_k$ computed on the sample, with $\sigma_1,\dots,\sigma_n$ i.i.d.\ Rademacher variables.\\
\textit{Proof in Appendix~\ref{prop:B1}.}
\end{proposition}

This proposition provides a formal bound on the best achievable risk. For simplicity, we assume $\hat{\mathcal{R}}^{k}_{n}(Q) = \hat{\mathcal{R}}^{k}(Q)$ throughout the rest of the paper to reduce notational complexity. Pre-trained models on various benchmark datasets are widely available in the literature. We argue that such models can be leveraged to estimate both the optimal regression order and the best achievable risk on the dataset they were trained on.

\begin{corollary}\label{prop:oracle}
Suppose we have access to a trained model $Q_k \in \mathcal{H}_k$ for each regression order $k = 1, \dots, M$. Then:
\[
  \begin{array}{r@{\quad}c}
    (\mathrm{i}) & \displaystyle\hat{\mathcal{R}}^{\infty}(Q^{*})
      = \min_{1\le k\le M}\{\hat{\mathcal{R}}^{k}(Q_k) - \Lambda(k)\},\\[6pt]
    (\mathrm{ii}) & \displaystyle k^*
      \leq \arg\min_{1\le k\le M}\{\hat{\mathcal{R}}^{k}(Q_k) - \Lambda(k)\}.
  \end{array}
\]
$(\mathrm{i})$ provides a way to estimate the minimal achievable risk. $(\mathrm{ii})$ allows one to infer the optimal regression order from this same model collection.\\
\textit{Proof in Appendix~\ref{cor:B2}.}
\end{corollary}

If model performance plateaus beyond a certain context length \( k \)—that is, if the loss ceases to improve while \( \Lambda(k) \) remains strictly positive—then the oracle reveals a lower achievable loss than what the model currently attains. This discrepancy indicates the presence of residual predictive structure that the model fails to capture, thereby offering a principled target for improvement. Oracle-based analysis thus provides a framework for assessing whether a model fully exploits the predictive information present in the data. Unlike \texttt{EvoRate}~\cite{zeng2025towards} or the forecastability measure~\cite{goerg2013forecastable}, our estimator is model-dependent, as it relies on the hypothesis class \( \mathcal{H}_k \). While this introduces sensitivity to the model architecture, it also enables direct comparisons between the estimated oracle risk and the empirical performance of candidate models, offering a more actionable diagnostic.

Finally, under a Mean Squared Error (MSE) loss and assuming the predictive distribution is multivariate Gaussian, it holds that $\mathcal{L}_{\mathrm{mle}} = \mathcal{L}_{\mathrm{mse}} + \text{const}$~\cite{zeng2025towards}. This equivalence ensures that our framework remains applicable in standard Gaussian regression settings, thereby enhancing its practical utility.

\section{Experiment}
\label{sec:Experiment}

Throughout this section, we use empirical estimators of the predictive information $\mathbf{I}_{\text{pred}}$, following the procedure outlined in Appendix~\ref{appendix:partB}.
For clarity, we will refer to our estimate simply as $\mathbf{I}_{\text{pred}}$ in the remainder of this section. In the first part, we explain the estimation procedure for $\mathbf{I}_{\text{pred}}$, and then use it to derive the estimated learning curve $\hat{\Lambda}$ and the optimal risk $\hat{\mathcal{R}}^{\infty}(Q^*)$.

\subsection{Estimation of $\mathbf{I}_{\text{pred}}$ on a Gaussian process}

We consider the process \( \{X_t\}_{t-k+1}^{t+k'} \), a \( d \)-dimensional Gaussian process with independent and identically distributed (i.i.d.) components. In this setting, it can be shown that
\begin{equation}
\label{eq:gaussien_process_evoRate_formula}
\mathbf{I}_{\text{pred}}(k,k') = \mathbf{I}(\mathbf{X}_{t-k+1}^{t}; \mathbf{X}_{t+1}^{t+k'}) = \frac{d}{2} \ln \left( \frac{|\Sigma_1^{(1)}||\Sigma_2^{(1)}|}{|\Sigma^{(1)}|} \right)
\end{equation}
where \( \Sigma_1 \) denotes the covariance matrix of \( \mathbf{X}_{t-k+1}^{t} \), \( \Sigma_2 \) the covariance of \( \mathbf{X}_{t+1}^{t+k'} \), and \( \Sigma \) the joint covariance matrix of the full sequence \( \{X_t\}_{t-k+1}^{t+k'} \).
\footnote{The superscript $(1)$ indicates that, due to the i.i.d. nature of the dimensions, the computation can be performed on the first dimension only and scaled by $d$.}

Figure~\ref{fig:estimation_evoPred} illustrates the estimation results of $\mathbf{I}_{\text{pred}}(k, k^\prime)$ across various input dimensions and kernel types detailed in Appendix~\ref{sub_section:synthetic_data_kernel}. For this experiment, we fixed $k=5$ and $k'=10$. $\mathbf{I}_{\text{pred}}$--\texttt{True} denotes the theoretical value computed from Equation~\ref{eq:gaussien_process_evoRate_formula}. Additional combinations of $(k, k')$ and their corresponding estimation results are provided in Appendix~\ref{sub_section:additionnal_combinations}.\\


\begin{figure}[H]
    \centering
    \includegraphics[width=1\linewidth]{figures/summary_k_5_k_prime_10.png} 
    \caption{
        Estimation of $\mathbf{I}_{\text{pred}}(k, k')$ using various neural-based methods. 
        Color encodes the estimation bias: blue regions indicate negative bias (underestimation), 
        while the intensity reflects the magnitude of this bias.
    }
    \label{fig:estimation_evoPred}
\end{figure}

While all methods provide relatively accurate estimates in low-dimensional settings (dimensionality $d \leq 20$), their performance deteriorates as the input space becomes more complex. Notably, methods like $\hat{\mathbf{I}}_{\text{pred}}$\text{-}\texttt{SMILE} and $\hat{\mathbf{I}}_{\text{pred}}$\text{-}\texttt{NWJ} exhibited consistent underestimation, in structured, high-dimensional regimes (e.g., Periodic and RQ kernels). Although these estimators inherently introduce variance and may lead to some instability in the results, refining the estimation of $\mathbf{I}_{\text{pred}}$ lies beyond the scope of this work and would warrant a dedicated investigation.
\subsection{Autoregressive Process}

We demonstrate the relevance of using predictive information $\mathbf{I}_{\text{pred}}$ to approximate the universal learning curve $\Lambda(k)$, as stated in Proposition~\ref{prop:Approximation_of_the_Universal_Learning_Curve}. This experiment also serves to validate the effectiveness of $\Lambda(k)$ in identifying the true memory length of a Markovian process, in accordance with Proposition~\ref{prop:markov_combined}. To this end, we simulate a stationary vector autoregressive process $\{X_t\}_{t=0}^{N-1} \subset \mathbb{R}^3$ of order $p \in \{5, 10\}$. The initial states $X_0, \dots, X_{p-1}$ are drawn independently from a standard multivariate normal distribution $\mathcal{N}(0, I_3)$. For $t \geq p$, the process evolves according to:
\begin{equation}
    X_t = \frac{\rho}{p} \sum_{j=t-p}^{t-1} X_j + \sqrt{1 - \rho^2} \, \epsilon_t,
    \label{eq:AR6}
\end{equation}
where $\epsilon_t \sim \mathcal{N}(0, I_3)$ and the parameter $\rho \in (0,1)$ controls the strength of temporal dependence.

Figure~\ref{fig:autoregressive_process_order_estimation} compares the estimated learning curve $\hat{\Lambda}(k)$—obtained via the estimator of predictive information from Proposition~\ref{prop:Approximation_of_the_Universal_Learning_Curve}—with a reference curve $\tilde{\Lambda}(k)$ computed from the known data-generating distribution. Although this theoretical curve is not available in real-world scenarios, it serves here as a useful benchmark. Since the full distribution is known, we can accurately approximate the entropy and thereby the true learning curve (see Appendix~\ref{sub_section:learning_curve} for derivation details).

The results confirm that $\hat{\Lambda}(k)$ closely tracks the theoretical curve and successfully identifies the correct model order $k = p$. This supports both the statistical consistency of the estimator and the practical usefulness of the learning curve for model selection. Despite minor fluctuations due to estimation variance, the method exhibits a sharp transition at the correct order $k = p$, underscoring its robustness and precision in capturing the underlying temporal structure.

\begin{figure}[H]
    \centering
    \includegraphics[width=0.8\textwidth]{figures/learning_curves.png}
    \caption{Learning curves $\Lambda(k)$ for AR processes with orders $p=5$ and $p=10$.}
    \label{fig:autoregressive_process_order_estimation}
\end{figure}

\begin{remark}
The ``Critical zone'' designates the range where the universal learning curve $\Lambda(k)$ falls below a small threshold (here set at $0.02$). Because estimation does not yield exact zeros due to inherent variance, it is difficult to pinpoint the exact context length at which the curve vanishes. Introducing such a cutoff therefore provides a practical criterion: values of $\Lambda(k)$ below this level are considered effectively zero. While the threshold is somewhat arbitrary, it is chosen small relative to the typical magnitude of $\Lambda(k)$, yet large enough to remain robust to estimation noise.
\end{remark}


\subsection{Estimating $\mathcal{R}^{\infty}(Q^*)$ for Ising Spin Sequences}
\label{sec:ising}
Let $\mathbf{X}_t^T = \{X_u\}_{u=t}^T \in \mathcal{X}^{T - t + 1}$ denote a sequence of binary variables where $\mathcal{X} = \{-1, +1\}$. We generate a sequence of length $T = 10{,}000{,}000$ starting with $X_0 \sim \text{Uniform}\{-1, +1\}$. The sequence is divided into blocks of size $M$. For each block, we sample $J \sim \mathcal{N}(0, 1)$ and evolve the sequence according to:
\begin{equation}
\label{eq:evolution_of_ising_chain}
P(X_i = +1 \mid X_{i-1}, J) = \frac{\exp(J X_{i-1})}{\exp(J X_{i-1}) + \exp(-J X_{i-1})}.
\end{equation}
This yields a piecewise-stationary Markov chain resembling a blockwise-random Ising process. We train both an MLP and an LSTM to predict $X_{t+1}$ from the $k$ past observations. Since an efficient predictor only requires a single parameter (logistic regression suffices), we fix $\dim\Theta = 1$, and use cross-entropy as the loss. For each block size \( M \in \{10{,}000,\ 100{,}000,\ 1{,}000{,}000, 10{,}000{,}000\} \), we evaluate, for \( 1 \leq k \leq n \) (with \( n = 19 \)), \texttt{EvoRate}(k) , along with our model-dependent estimates of the minimal achievable risk and the minimal loss empirically attained by LSTM and MLP models.

% \begin{table}[ht]
%     \centering
%     \begin{tabular}{c c c c c c}
%         \toprule
%          $M$ & $\texttt{EvoRate}(10)$ & $\hat{\mathcal{R}}^{\infty}_{\text{lstm}}(Q^*)$ & $\hat{\mathcal{R}}^{\infty}_{\text{mlp}}(Q^*)$ & $\min\limits_{1\leq k\leq n} 
%          \hat{\mathcal{R}}^{k}_{\text{lstm}}(Q)$ & $\min\limits_{1\leq k\leq n} \hat{\mathcal{R}}^{k}_{\text{mlp}}(Q)$ \\
%         \midrule
%         10{,}000 & 0.2758 & 0.3267& 0.3262& 0.4853 & 0.4872 \\
%         100{,}000 & 0.2861 & 0.3292 & 0.3329 & 0.4679 & 0.4668 \\
%         1{,}000{,}000  & 0.3269 & 0.3044 & 0.2715 & 0.3798 & 0.3660 \\
%         10{,}000{,}000  & 0.4760 & 0.0100 & 0.0860 & 0.0133 & 0.0903 \\
%         \bottomrule
%     \end{tabular}
%     \caption{Estimated minimal achievable risk and estimated minimal reached risk versus \texttt{EvoRate}.}
%     \label{tab:loss_comparison}
% \end{table}

\begin{table}[ht]
    \centering
    \begin{tabular}{c c c c c c}
        \toprule
         $M$ & $\texttt{EvoRate}(10)$ & $\hat{\mathcal{R}}^{\infty}_{\text{lstm}}(Q^*)$ & $\hat{\mathcal{R}}^{\infty}_{\text{mlp}}(Q^*)$ & $\min\limits_{1\leq k\leq n} 
         \hat{\mathcal{R}}^{k}_{\text{lstm}}(Q)$ & $\min\limits_{1\leq k\leq n} \hat{\mathcal{R}}^{k}_{\text{mlp}}(Q)$ \\
        \midrule
        10{,}000 & 0.2758 & 0.3724& 0.3719& 0.4853 & 0.4872 \\
        100{,}000 & 0.2861 & 0.3684 & 0.3720 & 0.4679 & 0.4668 \\
        1{,}000{,}000  & 0.3269 & 0.3569 & 0.3390 & 0.3798 & 0.3660 \\
        10{,}000{,}000  & 0.4760 & 0.0697 & 0.0867 & 0.0733 & 0.0903 \\
        \bottomrule
    \end{tabular}
    \caption{Estimated minimal achievable risk and estimated minimal reached risk versus \texttt{EvoRate}.}
    \label{tab:loss_comparison}
\end{table}
% \textit{See Appendix~\ref{sec:partC} for further details, statistical significance and results.}

As $M$ increases, the data becomes less complex: for $M = 10{,}000{,}000$, the coupling $J$ remains fixed, and the process reduces to a first-order Markov chain. Consequently, \texttt{EvoRate} increases, reflecting stronger underlying structure, and models achieve lower prediction losses. However, while \texttt{EvoRate} signals learnability, it does not provide a quantitative target, unlike our estimator $\hat{\mathcal{R}}^{\infty}(Q^*)$. Despite being model-dependent, the oracle estimates $\hat{\mathcal{R}}^{\infty}(Q^*)$ remain consistent across LSTM and MLP, with only minor differences relative to their scale. It also aligns qualitatively with \texttt{EvoRate}: lower estimated risk corresponds to more predictable sequences. Importantly, it enables direct comparison with actual model performance. For example, when $M = 10{,}000$, the ratio $\hat{\mathcal{R}}^{k}(Q) / \hat{\mathcal{R}}^{\infty}(Q^*) \approx 1.3$ indicates suboptimal predictions, likely due to high non-stationarity. As $M$ grows, this ratio approaches $1$, confirming improved model adequacy.
Lastly, negative $\Lambda(k)$ estimates for $M = 10{,}000{,}000$ stem from instability (\textit{see Table \ref{tab:universal_learning_4}}) in $\hat{\Lambda}(k)$ when $k \gg p$ for a true order-$p$ Markov process (\textit{see Figure~\ref{fig:autoregressive_process_order_estimation}}). Refining this estimator is necessary to avoid misinterpretation in low-complexity settings.

\subsection*{Additional Remarks on Result Interpretability}
We estimate the parameter dimension \( \dim(\Theta) \) using Corollary \ref{cor:universal_learning_curve_parametric}, with the true value being \( \mathbf{1} \). For \( k = 10 \), the estimator yields \( \hat{p} = 2 \times k \times \hat{\Lambda}(k) = 0.9580 \) when \( M = 10{,}000 \). Comparable values are obtained for other sample sizes \( M \), indicating the consistency of the procedure in the parametric regime. In contrast, when \( M = 10{,}000{,}000 \), which corresponds to the chain length, the process effectively becomes Markovian and thus departs from the parametric setting. As theoretically expected, the estimator then approaches zero; for example, at \( k = 10 \) we obtain \( \hat{p} = 0.0720 \).

We can also evaluate the optimal regression orders \( k^* \) for each model, which remain consistent across both, evolving together and staying within the same order of magnitude as \( M \) increases.

\begin{table}[ht]
    \centering
    \begin{tabular}{c c c}
        \toprule
        \( M \) & \( k^*_{\text{LSTM}} \) & \( k^*_{\text{MLP}} \) \\
        \midrule
        10{,}000     & 1 & 1 \\
        100{,}000    & 1 & 1 \\
        1{,}000{,}000  & 18 & 16 \\
        10{,}000{,}000 & 10 & 9 \\
        \bottomrule
    \end{tabular}
    \caption{Estimated optimal regression orders.}
    \label{tab:optimal_regression_orders}
\end{table}

\section{Conclusion}

% This paper introduces a unified framework that addresses two fundamental questions in sequential learning: \textbf{(i)} what is the minimal achievable risk for a predictor modeling sequential data, and \textbf{(ii)} can poor performance be attributed to model limitations or to the inherent unpredictability of the data? By generalizing the \texttt{EvoRate} metric through predictive information, we establish theoretical guarantees showing that temporal patterns fundamentally constrain learnability in sequential models.
% Our analysis provides explicit characterizations of the information-theoretic learning curve in different regimes. For Markov processes, we derive closed-form relationships between the learning curve and the Markov order. In parametric settings, we show that the learning curve decays proportionally to the inverse of the context length, with a coefficient determined by the dimensionality of the parameter space.
% Through experiments on synthetic datasets, we validated our framework for model selection and structure discovery. It accurately estimates the optimal achievable risk and distinguishes whether prediction performance is limited by model capacity or by the intrinsic complexity of the data.

\label{sec:conclusion}

This work addresses two fundamental questions in sequential modeling: 
\textbf{(i)} what is the minimal achievable risk when modeling sequential data, and 
\textbf{(ii)} is poor predictive performance due to model limitations or to the intrinsic unpredictability of the data?

To answer these questions, we introduced a unified information-theoretic framework centered around the learning curve 
$\Lambda(k)$, which quantifies the gain in predictive accuracy when extending the context from $k$ to $k+1$. 
This quantity provides a principled way to connect statistical dependencies in the data to achievable predictive performance. Building on this foundation, we proposed a practical estimator of the minimal achievable risk, $\hat{\mathcal{R}}^{\infty}(Q^{*})$ which bounds from below the performance of any predictor operating on the same data. This estimator enables a direct diagnostic test for model adequacy: if the empirical risk $\hat{\mathcal{R}}^{k}(Q)$ approaches $\hat{\mathcal{R}}^{\infty}(Q^{*})$, then the observed performance is close to the intrinsic unpredictability of the process—suggesting that increasing model capacity or context length is unlikely to yield further improvements. 
Conversely, a significant gap between these two quantities reveals that the model underfits the data, pointing to unexploited temporal dependencies.

Through theoretical analysis, we demonstrated that $\Lambda(k)$ admits explicit asymptotic forms in both parametric and Markov regimes, linking the decay of the learning curve to intrinsic properties such as the parameter dimensionality or the Markov order. 
Empirical validation on controlled synthetic datasets confirmed these predictions: our estimator consistently recovered the true minimal risk, accurately distinguished between parametric and non-parametric regimes, and correctly identified whether performance limitations arose from model capacity or from the intrinsic randomness of the process.

% In summary, this work provides a unified information-theoretic framework that quantifies the fundamental limits of sequential prediction and distinguishes model insufficiency from intrinsic data unpredictability. By linking predictive information to empirical risks, it offers a principled way to estimate the minimal achievable risk and to assess whether observed performance reflects model capacity or data complexity.

\newpage

\bibliography{references}  

%%%%%%%%%%%%%%%%%%%%%%%%%%%%%%%%%%%%%%%%%%%%%%%%%%%%%%%%%%%%


\newpage

\appendix
\section*{Appendix}
\addcontentsline{toc}{section}{Appendix}

\startcontents[appendix]
\printcontents[appendix]{l}{1}{\setcounter{tocdepth}{2}}

\clearpage

\section{Proofs}
\label{appendix:partA}
\subsection{On the main properties of $\mathbf{I}_{\text{pred}}$}


\begin{proposition}[Elementary properties of the predictive mutual information]
\label{prop:basic_Ipred_properties}
For integers $k\!\ge 1$ and $k'\!\ge 1$, let
\[
\mathbf I_{\mathrm{pred}}(k,k')
\;=\;
\mathbf I\!\bigl(\mathbf X_{t-k+1}^{t}\,;\, \mathbf X_{t+1}^{t+k'}\bigr)
\]
denote the mutual information between a past block of length $k$ and a future
block of length $k'$.  
Assuming that the process $(X_t)_{t\in\mathbb Z}$ is
stationary (and therefore time–translation invariant),  
$\mathbf I_{\mathrm{pred}}$ satisfies:

\begin{enumerate}
    \item \textbf{Non-negativity.}\;
          $\mathbf I_{\mathrm{pred}}(k,k') \ge 0$.

    \item \textbf{Symmetry.}\;
          $\mathbf I_{\mathrm{pred}}(k,k')
           = \mathbf I\!\bigl(\mathbf X_{t+1}^{t+k'}\,;\, \mathbf X_{t-k+1}^{t}\bigr)$.

    \item \textbf{Monotonicity.}
          \begin{enumerate}
              \item (Increasing past) For fixed $k'$,  
                    $\mathbf I_{\mathrm{pred}}(k+1,k')\;\ge\;
                     \mathbf I_{\mathrm{pred}}(k,k')$.
              \item (Increasing future) For fixed $k$,  
                    $\mathbf I_{\mathrm{pred}}(k,k'+1)\;\ge\;
                     \mathbf I_{\mathrm{pred}}(k,k')$.
          \end{enumerate}

    \item \textbf{Chain-rule decomposition.}\;
          For any split $k_1+k_2=k'$,
          \[
            \mathbf I_{\mathrm{pred}}(k,k_1+k_2)
            = \mathbf I_{\mathrm{pred}}(k,k_1)
            + \mathbf I\!\bigl(
                 \mathbf X_{t-k+1}^{t}\,;\,
                 \mathbf X_{t+k_1+1}^{t+k_1+k_2}
                 \,\bigm|\,
                 \mathbf X_{t+1}^{t+k_1}
              \bigr),
          \]
          and an analogous identity holds when the past block is partitioned.

    \item \textbf{Data-processing inequality.}\;
          For any measurable maps $f$ and $g$,
          \[
            \mathbf I\!\bigl(
               f(\mathbf X_{t-k+1}^{t})\,;\,
               g(\mathbf X_{t+1}^{t+k'})
            \bigr)
            \;\le\;
            \mathbf I_{\mathrm{pred}}(k,k').
          \]

    \item \textbf{Convergence to excess entropy.}\;
          Let $E$ denote the excess (or predictive) entropy of the process:
          \(
             E := \mathbf I\!\bigl(\mathbf X_{-\infty}^{t}\,;\,
                                   \mathbf X_{t+1}^{\infty}\bigr).
          \)
          Then  
          $\displaystyle\lim_{k,k'\to\infty}\mathbf I_{\mathrm{pred}}(k,k') = E$.
\end{enumerate}
\end{proposition}

\begin{proof}
All items follow directly from classical properties of Shannon’s
mutual information (see, e.g., \cite[Chs.~2–3]{cover1999elements}).

\begin{itemize}
    \item \emph{(1)–(2)} are immediate from non-negativity and symmetry of~$I(X;Y)$.

    \item \emph{(3)}  
          Extending the conditioning set can only reduce conditional entropy,
          hence cannot decrease $I(X;Y)$; apply this with  
          $X=\mathbf X_{t-k}^{t}$ (resp.\ $\mathbf X_{t-k+1}^{t}$) and  
          $Y=\mathbf X_{t+1}^{t+k'}$ (resp.\ $\mathbf X_{t+1}^{t+k'+1}$).

    \item \emph{(4)} is the chain rule
          $I(X;YZ)=I(X;Y)+I(X;Z\mid Y)$
          applied to suitably chosen blocks.

    \item \emph{(5)} is the data-processing inequality:
          applying measurable maps cannot increase mutual information.

    \item \emph{(6)}  
          For fixed $k$, $\mathbf I_{\mathrm{pred}}(k,k')$ is non-decreasing and
          bounded above by $E$; likewise in $k$.  
          Monotone convergence plus stationarity gives the limit $E$.
\end{itemize}
\end{proof}


\begin{proposition}[Convergence of $\mathbf I_{\mathrm{pred}}$ when the future window grows]
\label{prop:Ipred_to_subextensive}
Let $(X_t)_{t\in\mathbb Z}$ be a stationary process with finite
entropy rate $h_0:=\lim_{n\to\infty}H(X_1^n)/n$.
Write the block entropy as
\[
   H(n)=n\,h_0+H_1(n),
   \qquad n\ge 1,
\]
where the \emph{sub-extensive term} satisfies $H_1(n)=o(n)$.
Then, for every fixed $k\ge 1$,
\[
   \lim_{k'\to\infty}
   \mathbf I_{\mathrm{pred}}(k,k')
   \;=\;
   H_1(k).
\]
\end{proposition}

\begin{proof}
For any $k,k'\ge 1$ the predictive mutual information can be written as  
\[
   \mathbf I_{\mathrm{pred}}(k,k')
   \;=\;
   H(k)+H(k')-H(k+k').
\]
Insert the decomposition $H(n)=nh_0+H_1(n)$:
\[
\mathbf I_{\mathrm{pred}}(k,k')
  = \bigl[kh_0+H_1(k)\bigr]
    +\bigl[k'h_0+H_1(k')\bigr]
    -\bigl[(k+k')h_0+H_1(k+k')\bigr]
  = H_1(k)+H_1(k')-H_1(k+k').
\]
Now let $k'\to\infty$ while keeping $k$ fixed.
Because $H_1(n)=o(n)$, the difference $H_1(k')-H_1(k+k')$ vanishes:
\[
   |H_1(k')-H_1(k+k')|
   \;\le\;
   o(k')+o(k+k')=o(k'),
\]
hence tends to $0$.  
Therefore
\(
  \displaystyle
  \lim_{k'\to\infty}\mathbf I_{\mathrm{pred}}(k,k') = H_1(k).
\)
\end{proof}



\begin{proposition}[Asymptotic equivalence between $H(k)$ and $\mathbf I_{\mathrm{pred}}$]
\label{prop:Ipred_vs_entropyblock}
Assume Hypothesis $\mathbf{(H_0)}$ (stationarity with finite entropy
rate $l_0$ and sub-extensive remainder $H_1$).
Fix a non-decreasing sequence $k'\!=k'(k)$ such that $k'\!\ge k$ and
$k'\!=\mathcal O(k)$ when $k\to\infty$.
Then
\[
   H\!\bigl(\mathbf X_{t-k+1}^{t}\bigr)
   \;=\;
   k\,l_{0}
   +\mathbf I_{\mathrm{pred}}\!\bigl(k,k'\bigr)
   +o(k),
   \qquad k\to\infty,
\]
hence in particular
$\displaystyle\mathbf I_{\mathrm{pred}}(k,k')/k \;\longrightarrow\;0$.
\end{proposition}

\begin{proof}
Write the block entropy as $H(n)=n\,l_{0}+H_{1}(n)$ with
$H_{1}(n)=o(n)$.
Using the identity
$\mathbf I_{\mathrm{pred}}(k,k')
   =H(k)+H(k')-H(k+k')$
and the decomposition above gives
\[
   \mathbf I_{\mathrm{pred}}(k,k')
   = H_{1}(k)+H_{1}(k')-H_{1}(k+k').
\]
Because $k'\ge k$ and $k'=\mathcal O(k)$,
sub-extensiveness yields
$H_{1}(k')-H_{1}(k+k') = o(k)$.
Therefore
\[
   \mathbf I_{\mathrm{pred}}(k,k')
   = H_{1}(k) + o(k).
\]
Insert this into $H(k)=k\,l_{0}+H_{1}(k)$ to obtain
\(
   H(k)=k\,l_{0}+\mathbf I_{\mathrm{pred}}(k,k')+o(k).
\)
Finally,
$\mathbf I_{\mathrm{pred}}(k,k')/k
   = H_{1}(k)/k + o(1)\xrightarrow{k\to\infty}0$
because $H_{1}(k)=o(k)$.
\end{proof}

\begin{proposition}[Exact link between \texttt{EvoRate} and $\mathbf I_{\mathrm{pred}}$]
\label{prop:EvoRate_vs_Ipred}
Assume the process is stationary and time–translation invariant.  
For any integers $k\!\ge1$ and $k'\!\ge1$,
\[
   \mathbf I_{\mathrm{pred}}(k,k')
   \;=\;
   \texttt{EvoRate}(k)
   \;+\;\bigl[\,H_{1}(k')-H_{1}(1)\bigr]
   \;+\;\bigl[\,H_{1}(k+1)-H_{1}(k+k')\bigr],
\]
where the block entropy is decomposed as
$H(n)=n\,l_{0}+H_{1}(n)$ with a sub-extensive part $H_{1}(n)=o(n)$.
In particular, for a single-step future window ($k'=1$) we recover
\[
   \mathbf I_{\mathrm{pred}}(k,1)=\texttt{EvoRate}(k).
\]
\end{proposition}

\begin{proof}
By stationarity,
\(
   \mathbf I_{\mathrm{pred}}(k,k')
   = H_{1}(k')+H_{1}(k)-H_{1}(k+k').
\)
The evolutionary rate is defined as
\(
   \texttt{EvoRate}(k)=H_{1}(1)+H_{1}(k)-H_{1}(k+1).
\)
Solving the last equation for $H_{1}(k)$ gives  
$H_{1}(k)=\texttt{EvoRate}(k)-H_{1}(1)+H_{1}(k+1)$.
Substituting this into the expression of $\mathbf I_{\mathrm{pred}}$
yields
\[
   \mathbf I_{\mathrm{pred}}(k,k')
   =\texttt{EvoRate}(k)
    +\bigl[H_{1}(k')-H_{1}(1)\bigr]
    +\bigl[H_{1}(k+1)-H_{1}(k+k')\bigr],
\]
which is the desired identity.  
Setting $k'=1$ cancels both bracketed terms, proving the special case
$\mathbf I_{\mathrm{pred}}(k,1)=\texttt{EvoRate}(k)$.
\end{proof}




\subsection{Learning theory}

\begin{proposition}[Predictive information versus the cumulative learning curve]
\label{prop:Ipred_vs_Lambda_bounds_corrected}
Assume \textbf{(H\textsubscript{0})} (strict stationarity and a finite
entropy rate).  
For \(n\ge 1\) write the block entropy as \(H(n)=n\,\ell_{0}+H_{1}(n)\),
where the \emph{sub-extensive} term satisfies \(H_{1}(n)=o(n)\).
Define the order–\(k\) entropy rate
\(
  \ell(k)=H(k+1)-H(k)
\)
and the \emph{universal learning curve}
\(
  \Lambda(k)=\ell(k)-\ell_{0}
            =H_{1}(k+1)-H_{1}(k).
\)

Then, for all integers \(k\ge 1\) and \(k'\ge k\),
\begin{equation}\label{eq:Ipred_bounds_final}
   \sum_{i=1}^{k}\Lambda(i)\;-\;H_{1}(k)
   \;\;\le\;\;
   \mathbf I_{\mathrm{pred}}(k,k')
   \;\;\le\;\;
   \sum_{i=1}^{k}\Lambda(i).
\end{equation}
Moreover
\begin{equation}\label{eq:kLambda_limit}
   k\,\Lambda(k)\;\xrightarrow[k\to\infty]{}\;0,
\end{equation}
so the width of the sandwich in \eqref{eq:Ipred_bounds_final} is
\(H_{1}(k)=o(k)\).
\end{proposition}

\begin{proof}
From the decomposition \(H(n)=n\ell_{0}+H_{1}(n)\) we have
\(
  \Lambda(k)=H_{1}(k+1)-H_{1}(k)
\)
and
\(
  S_{k}:=\sum_{i=1}^{k}\Lambda(i)=H_{1}(k+1)-H_{1}(1).
\)

For any \(k'\ge k\),
\[
\mathbf I_{\mathrm{pred}}(k,k')
  =H(k)+H(k')-H(k+k')
  =H_{1}(k)+H_{1}(k')-H_{1}(k+k').
\]
Hence
\begin{equation}\label{eq:aux1}
  \mathbf I_{\mathrm{pred}}(k,k')
  =S_{k}-\bigl[H_{1}(k+k')-H_{1}(k')\bigr].
\end{equation}

Because \(H_{1}\) is non-decreasing,
\(H_{1}(k+k')\ge H_{1}(k')\); the bracket in
\eqref{eq:aux1} is therefore non-negative and we obtain the upper bound
\(
  \mathbf I_{\mathrm{pred}}(k,k')\le S_{k}.
\)

Sub-additivity of \(H_{1}\) gives
\(H_{1}(k+k')\le H_{1}(k')+H_{1}(k)\), hence
\(H_{1}(k+k')-H_{1}(k')\le H_{1}(k)\).
Inserting this into \eqref{eq:aux1} yields the lower bound
\(
  \mathbf I_{\mathrm{pred}}(k,k')\ge S_{k}-H_{1}(k).
\)
Together these inequalities establish \eqref{eq:Ipred_bounds_final}.

Finally,
\(
  k\,\Lambda(k)
  =k\,\bigl[H_{1}(k+1)-H_{1}(k)\bigr]
  \le H_{1}(k+1)=o(k),
\)
which proves \eqref{eq:kLambda_limit}.
\end{proof}





\begin{proposition}[Predictive–information increment vs.\ universal learning curve]
\label{prop:learning_curve_approximation}
Fix \(k\ge 1\).  
Then, under hypothesis \textbf{(H\textsubscript{0})},
\[
   \mathbf I_{\mathrm{pred}}\!\bigl(k\!+\!1,k'\bigr)
   \;-\;
   \mathbf I_{\mathrm{pred}}\!\bigl(k,k'\bigr)
   \;\xrightarrow[k'\to\infty]{}\;
   \Lambda(k),
\]
where \(\Lambda(k)=\ell(k)-\ell_0\) is the universal learning curve and  
\(\ell(k)=H(k+1)-H(k)\) denotes the order-\(k\) entropy rate.
\end{proposition}

\begin{proof}

Write the block entropy in its extensive–plus–remainder form  
\(H(n)=n\,\ell_0+H_1(n)\) with \(H_1(n)=o(n)\).
Consequently
\[
   \ell(k)=H(k+1)-H(k)=H_1(k+1)-H_1(k),
   \qquad
   \Lambda(k)=\ell(k)-\ell_0
             =H_1(k+1)-H_1(k).
\tag{1}\label{eq:Lambda_H1}
\]

Proposition \ref{prop:Ipred_to_subextensive} (convergence to the
sub-extensive part) gives, for every fixed \(m\),
\[
   \lim_{k'\to\infty}\mathbf I_{\mathrm{pred}}(m,k')
   \;=\;
   H_1(m).
\tag{2}\label{eq:Ipred_to_H1}
\]

Applying \eqref{eq:Ipred_to_H1} with \(m=k\) and \(m=k+1\) we obtain
\[
   \lim_{k'\to\infty}
   \bigl[
      \mathbf I_{\mathrm{pred}}(k+1,k')
      -\mathbf I_{\mathrm{pred}}(k,k')
   \bigr]
   =H_1(k+1)-H_1(k).
\]
Combined with \eqref{eq:Lambda_H1}, this equals \(\Lambda(k)\),
establishing the claimed convergence.
\end{proof}


\subsection{Asymptotic behavior of $\mathbf{I}_{\text{pred}}$}

\subsubsection{Markovian case}
We believe that the main limitations of \texttt{EvoRate} stem from the fact that its limiting values are unclear. It is difficult to conclude about its limits. 
In contrast, the study of $\mathbf{I}_{\text{pred}}$
and its limiting values can provide meaningful insights and may help to uncover underlying patterns in sequential data. This is particularly true in the case of Markovian processes:

\begin{proposition}\label{prop:value_of_evoPred_for_markovian_process}
Assume that $X$ is a Markov process of order $m$ and that $k'\geq k \geq m$ so that the relevant past information is contained in $\mathbf{X}_{t-m+1}^t$, then 
$$\mathbf{I}_{\text{pred}}\left(k, k^{\prime}\right)=I\left(\mathbf{X}_{t-k+1}^t, X_{t+1}^{t+k^{\prime}}\right)=\mathbb{E}_{\mathbf{X}_{t-m+1}^{t+m}}\left[\ln \frac{P\left(X_{t+1}^{t+m} \mid \mathbf{X}_{t-m+1}^t\right)}{P\left(X_{t+1}^{t+m}\right)}\right].$$
In particular for a order-1 Markov process, $\mathbf{I}_{\text{pred}}\left(k, k^{\prime}\right) = \texttt{EvoRate}(1)$.
\end{proposition}
\begin{proof}

We start from the definition 
$$
\mathbf{I}_{\text{pred}}\left(k, k^{\prime}\right)=I\left(\mathbf{X}_{t-k+1}^t, X_{t+1}^{t+k^{\prime}}\right)=\mathbb{E}_{\mathbf{X}_{t-k+1}^{t+k^{\prime}}}\left[\ln \frac{P\left(X_{t+1}^{t+k^{\prime}} \mid \mathbf{X}_{t-k+1}^t\right)}{P\left(X_{t+1}^{t+k^{\prime}}\right)}\right]
$$

Assume that $X$ is a Markov process of order $m$ and that $k \geq m$ so that the relevant past information is contained in $\mathbf{X}_{t-m+1}^t$. We start by decomposing the unconditionnal probability (the denominator)
$$
P\left(X_{t+1}^{t+k^{\prime}}\right)=\prod_{j=1}^{k^{\prime}} P\left(X_{t+j} \mid X_{t+1}^{t+j-1}\right)
$$

Since the process is Markov of order $m$, the simplification by the Markov property holds only when there are at least $m$ prior observations in the sequence $X_{t+1}^{t+j-1}$. For $j=1, \ldots, m$, the conditional probability $P\left(X_{t+j} \mid X_{t+1}^{t+j-1}\right)$ remains as is because $X_{t+1}^{t+j-1}$ contains fewer than $m$ observations (with the understanding that, by convention, for $j=1$ we have $P\left(X_{t+1} \mid \emptyset\right)=P\left(X_{t+1}\right)$ ). For $j \geq m+1$, the Markov property yields $P\left(X_{t+j} \mid X_{t+1}^{t+j-1}\right)=P\left(X_{t+j} \mid X_{t+j-m}^{t+j-1}\right) .$


Thus, the full decomposition is
\begin{equation}\label{eq:denominator_term_simplification}
P\left(X_{t+1}^{t+k^{\prime}}\right)=\underbrace{\prod_{j=1}^m P\left(X_{t+j} \mid X_{t+1}^{t+j-1}\right)}_{\text {non-simplified terms }} \times \underbrace{\prod_{j=m+1}^{k^{\prime}} P\left(X_{t+j} \mid X_{t+j-m}^{t+j-1}\right)}_{\text {Markov terms }} .
\end{equation}

Now, consider the numerator $P\left(X_{t+1}^{t+k^{\prime}} \mid \mathbf{X}_{t-k+1}^t\right)$. Since $k \geq m$, the available past $\mathbf{X}_{t-k+1}^t$ contains at least the last $m$ values, i.e., $X_{t-m+1}^t$. We directly have by similar arguments, 
\begin{equation}\label{eq:numerateur_term_simplification}
P\left(X_{t+1}^{t+k^{\prime}} \mid \mathbf{X}_{t-k+1}^t\right)=\prod_{j=1}^{k^{\prime}} P\left(X_{t+j} \mid X_{t-m+1}^{t+j-1}\right)
\end{equation}


Using \ref{eq:denominator_term_simplification} and \ref{eq:numerateur_term_simplification} the factors for $j \geq m+1$ in the numerator and the denominator match and cancel each other in the ratio. In other words, the difference between the conditional probability $P\left(X_{t+1}^{t+k^{\prime}} \mid \mathbf{X}_{t-k+1}^t\right)$ and the unconditional probability $P\left(X_{t+1}^{t+k^{\prime}}\right)$ is confined to the first $m$ factors. We therefore write

$$
\frac{P\left(X_{t+1}^{t+k^{\prime}} \mid \mathbf{X}_{t-k+1}^t\right)}{P\left(X_{t+1}^{t+k^{\prime}}\right)}=\frac{\prod_{j=1}^m P\left(X_{t+j} \mid X_{t-m+1}^{t+j-1}\right)}{\prod_{j=1}^m P\left(X_{t+j} \mid X_{t+1}^{t+j-1}\right)}
$$

We obtain, after simplification : 
$$
\frac{P\left(X_{t+1}^{t+k^{\prime}} \mid \mathbf{X}_{t-k+1}^t\right)}{P\left(X_{t+1}^{t+k^{\prime}}\right)}=\frac{P\left(X_{t+1}^{t+m} \mid X_{t-m+1}^t\right)}{P\left(X_{t+1}^{t+m}\right)}
$$
leading to 
$$\mathbf{I}_{\text{pred}}\left(k, k^{\prime}\right)=I\left(\mathbf{X}_{t-k+1}^t, X_{t+1}^{t+k^{\prime}}\right)=\mathbb{E}_{\mathbf{X}_{t-m+1}^{t+m}}\left[\ln \frac{P\left(X_{t+1}^{t+m} \mid \mathbf{X}_{t-m+1}^t\right)}{P\left(X_{t+1}^{t+m}\right)}\right]$$
\end{proof}



\begin{proposition}\label{proof:markovian_process_universal_learning_curve} Let's suppose $\mathbf{X}_{t}^{T}$ is a Markovian process of order $m$. Then 
$$\forall k, \qquad k\geq m \Rightarrow \Lambda(k) = 0,$$
\end{proposition}
\begin{proof}
From the definition of the universal learning curve, and noting that $\mathbf{I}_{\text{pred}}$ is constant for $k$ greater than $m$.
\end{proof}

\subsubsection{Predictive information for a parametrised stationary process}
\label{proof:parametrized_process_evoPred_limit}

Our goal is to derive the asymptotic expansion of the predictive mutual
information when the data–generating process belongs to a finite-dimensional
parametric family.  The argument extends \citet{bialek1999predictive}, who
treated the i.i.d.\ case, to the setting of \emph{dependent} sequences.

\paragraph{Setup and notation.}
Let $\{X_t\}_{t\in\mathbb Z}$ be a strictly stationary stochastic process
taking values in a Polish space $\mathcal X\subset\mathbb R^{d}$.  For
integers $k\ge 1$ and $k'\ge k$ denote
\[
  \mathbf X_{\mathrm{past}}
    := \mathbf X_{t-k+1}^{t}
    \;=\;(X_{t-k+1},\dots,X_{t}),
  \qquad
  \mathbf X_{\mathrm{fut}}
    := \mathbf X_{t+1}^{t+k'}
    \;=\;(X_{t+1},\dots,X_{t+k'}).
\]
Write $p(x_{1}^{n})$ for the joint density of the block
$\mathbf X_{1}^{n}:=(X_{1},\dots,X_{n})$ with respect to a reference measure
$\lambda^{d\otimes n}$ on $\mathcal X^{n}$.

\paragraph{Parametric model.}
Assume that there exists
\begin{itemize}
  \item an \emph{open} parameter set $\Theta\subset\mathbb R^{p}$ ($p<\infty$),
  \item a \emph{prior density} $\mathcal P\colon\Theta\to(0,\infty)$ of class
        $C^{1}$ on $\Theta$,
  \item a \emph{Kolmogorov-consistent} family of densities
        $\bigl\{Q^{(n)}_\theta\bigr\}_{\theta\in\Theta,n\ge 1}$ such that, for
        each $n\ge 1$,
        \[
            p(x_{1}^{n})
            \;=\;
            \int_{\Theta}
              Q^{(n)}_\theta(x_{1}^{n})\,
              \mathcal P(\theta)\,d\theta.
        \]
\end{itemize}
Consistency means that $(Q^{(n)}_\theta)_{n\ge 1}$ are the marginals of a
single probability law $Q_\theta$ on $\mathcal X^{\mathbb Z}$
\citep[Chap.\,8]{kallenberg2002foundations}.  The true parameter
$\bar\theta\in\Theta$ is the (unknown) value generating the observations.

\paragraph{Main regularity hypotheses.}

\begin{enumerate}

\item \label{H:ergodic} 
\textbf{Stationarity and geometric $\boldsymbol{\alpha}$-mixing.}\;
Under the true parameter $\bar\theta$, the process is strictly stationary and
ergodic, with Rosenblatt mixing coefficients satisfying
$\alpha_{\bar\theta}(n)\le C e^{-c n}$ for some constants $C,c>0$.

\item \label{H:ident} 
\textbf{$\boldsymbol{C^{3}}$ identifiability of the KL map.}\;
The function
$\theta\mapsto
  D_{\mathrm{KL}}\!\bigl(Q_\theta\,\Vert\,Q_{\bar\theta}\bigr)$
is three-times continuously differentiable on a neighbourhood of $\bar\theta$,
attains its unique minimum at $\bar\theta$, and has positive-definite Hessian
$\mathcal F$ at that point.

\item \label{H:entropy} 
\textbf{Finite entropy rate.}\;
The block entropy
$H(n):=H_{Q_{\bar\theta}}(\mathbf X_{1}^{n})$
satisfies the Shannon–McMillan property
$H(n)=n\,\ell_{0}+o(n)$ with $\ell_{0}<\infty$.

\end{enumerate}


\begin{theorem}[Asymptotics of the predictive mutual information]
\label{thm:PI}
Under assumptions \ref{H:ergodic}–\ref{H:entropy}, as $k\to\infty$ with
$k'\ge k$,
\begin{align}
  \mathbf I_{\mathrm{pred}}(k,k')
  &= I\!\bigl(\mathbf X_{\mathrm{past}};\mathbf X_{\mathrm{fut}}\bigr)
     \nonumber\\
  &= \frac{p}{2}\,\ln k
     + \frac12\ln\det\mathcal F
     - \frac{p}{2}\,\ln(2\pi)
     + \ln\mathcal P(\bar\theta)
     + \mathcal O\!\bigl(k^{-1}\bigr).
  \label{eq:PI-asympt}
\end{align}
Consequently the universal learning curve
$\Lambda(k):=\ell(k)-\ell_{0}$ obeys
\begin{equation}
  \Lambda(k)
  \;=\;
  \frac{p}{2k}
  -\frac{p}{4k^{2}}
  +\mathcal O\!\bigl(k^{-3}\bigr).
  \label{eq:Lambda-asympt}
\end{equation}
All $o(\cdot)$ and $\mathcal O(\cdot)$ symbols are uniform in
$\theta\in\mathcal N(\bar\theta)$.
\end{theorem}

\begin{proof}
Let $n:=k+k'$.  For $\theta\in\Theta$ set
$L_n(\theta):=
    \ln\!\bigl(Q^{(n)}_\theta(\mathbf X_{1}^{n})/
                Q^{(n)}_{\bar\theta}(\mathbf X_{1}^{n})\bigr)$.
Exponential $\alpha$-mixing together with a Bernstein-type blocking argument
\citep{dedecker2007weak,dedecker2003conditional} yields a \emph{uniform}
strong law of large numbers:
\begin{equation}
  \sup_{\theta\in\mathcal N(\bar\theta)}
  \Bigl|
      n^{-1}L_n(\theta)
      + D_{\mathrm{KL}}\!\bigl(Q_\theta\Vert Q_{\bar\theta}\bigr)
  \Bigr|
  \;\xrightarrow[n\to\infty]{\;Q_{\bar\theta}\textrm{-a.s.}}\;
  0.
  \label{eq:U-LLN}
\end{equation}
Assumption~\ref{H:ident} gives the quadratic expansion
$D_{\mathrm{KL}}(Q_\theta\Vert Q_{\bar\theta})
  =\tfrac12(\theta-\bar\theta)^{\!\top}\!\mathcal F(\theta-\bar\theta)
   +\mathcal O(\|\theta-\bar\theta\|^{3})$.
Plugging this into \eqref{eq:U-LLN} and integrating, we split the marginal
likelihood as
\[
  p(\mathbf X_{1}^{n})
  = Q^{(n)}_{\bar\theta}(\mathbf X_{1}^{n})
    \int_{\Theta} \exp\!\bigl(L_{n}(\theta)\bigr)\,
                 \mathcal P(\theta)\,d\theta
  =: Q^{(n)}_{\bar\theta}(\mathbf X_{1}^{n})\,\mathcal Z_n.
\]

\emph{Laplace expansion.}  Because
$D_{\mathrm{KL}}(\theta\Vert\bar\theta)\ge
  c\|\theta-\bar\theta\|^{2}$ in $\mathcal N(\bar\theta)$, decompose
$\Theta=B(\bar\theta,n^{-1/2+\delta})\cup\text{(complement)}$.  On the small
ball, $L_n(\theta)$ is dominated by the quadratic form
$-\frac n2(\theta-\bar\theta)^{\!\top}\!\mathcal F(\theta-\bar\theta)$ with a
cubic remainder uniformly $o(1)$.  Classical Laplace–Watson lemma
\citep[Chap.\,6]{bleistein2012asymptotic} yields
\begin{equation}
  \mathcal Z_n
  =(2\pi)^{p/2}\,n^{-p/2}\,
    \frac{\mathcal P(\bar\theta)}{\sqrt{\det\mathcal F}}\,
    \bigl\{1+\mathcal O(n^{-1})\bigr\}.
  \label{eq:Laplace}
\end{equation}
The contribution of the complement is $o\!\bigl(n^{-p/2}\bigr)$ because the
integrand decays like $\exp\!\bigl[-c n^{\delta}\bigr]$.

\emph{Block entropy decomposition.}
Taking the $-\ln$ and the $Q_{\bar\theta}$-expectation of
$Q^{(n)}_{\bar\theta}(\mathbf X_{1}^{n})\,\mathcal Z_n$
and invoking \ref{H:entropy} give
\begin{equation}
  H(n)
  = n\,\ell_{0}
    +\frac{p}{2}\ln n
    +\frac12\ln\det\mathcal F
    -\frac{p}{2}\ln(2\pi)
    +\ln\mathcal P(\bar\theta)
    +\mathcal O(n^{-1}).
  \label{eq:H(n)}
\end{equation}

\emph{Predictive information.}
By definition
$\mathbf I_{\mathrm{pred}}(k,k')=H(k)+H(k')-H(k+k')$.
Substituting \eqref{eq:H(n)} with $n=k,k',k+k'$ cancels the extensive
$n\ell_{0}$ parts; the remaining sub-extensive contributions yield
\eqref{eq:PI-asympt}.  Finally
\[
  \Lambda(k)=\ell(k)-\ell_{0}
            =H(k+1)-2H(k)+H(k-1)
            =\frac{p}{2k}-\frac{p}{4k^{2}}+\mathcal O(k^{-3}),
\]
where the last equality follows from a Taylor expansion
$\ln(k\!+\!1)-\ln k = k^{-1}-\tfrac12 k^{-2}+O(k^{-3})$.  Uniformity of the
remainders is ensured by the uniform Laplace estimate
\eqref{eq:Laplace}.  This completes the proof.
\end{proof}

\begin{remark}
The leading term $\tfrac p2\ln k$ coincides with the \emph{model-complexity
penalty} in Bayesian minimum description length
\citep{clarke1990information,barron1991minimum}.  The constant
$\ln\mathcal P(\bar\theta)-\tfrac p2\ln(2\pi)$ depends on the local prior
mass; it vanishes in the derivative $\Lambda(k)$ but is essential for the
absolute scale of $\mathbf I_{\mathrm{pred}}$.
\end{remark}


\subsection{Minimal achievable risk}

\begin{proposition}[Learning–curve surplus bound]\label{prop:B1}
Let\/ $(X_t)_{t\in\mathbb Z}$ be stationary and \emph{geometrically
$\alpha$–mixing}, i.e.\ $\alpha(n)\le C e^{-c n}$ for constants
$C,c>0$.  
For every regression order $k\in\mathbb N$, every predictor
$Q\in\mathcal H_k$, every sample size $n\ge1$ and every confidence
$\delta\in(0,1)$, with probability at least $1-\delta$
\begin{equation}\label{eq:B1-main}
   \mathcal R^{\infty}(Q^{*})
   \;\le\;
   \hat{\mathcal R}^{k}(Q)
   \;-\;\Lambda(k)
   \;+\;2\,\widehat{\Re}_n(\mathcal F_k)
   \;+\;
   3\,\frac{\ln(1/\delta)}{n},
\end{equation}
where $\widehat{\Re}_n(\mathcal F_k)$ is the empirical Rademacher
complexity of the loss class
\(
  \mathcal F_k=\bigl\{\,
      (\mathbf x,x')\mapsto -\ln Q(x'\mid\mathbf x)
      : Q\in\mathcal H_k
   \bigr\}.
\)
\end{proposition}

\begin{proof}
For any $Q\in\mathcal H_k$
\[
   \mathcal R^{k}(Q)
   \;=\;
   \mathcal R^{\infty}(Q^{*})
   \;+\;
   \bigl[\mathcal R^{k}(Q)-\mathcal R^{k}(Q^{*})\bigr]
   \;+\;\Lambda(k),
\]
because
$\Lambda(k)=\mathcal R^{k}(Q^{*})-\mathcal R^{\infty}(Q^{*})$.
Since the loss function $\ell(\cdot,\cdot)$ is non–negative,
$\mathcal R^{k}(Q)\ge\mathcal R^{k}(Q^{*})$, hence
\[
   \mathcal R^{k}(Q)
   \;\ge\;
   \mathcal R^{\infty}(Q^{*}) + \Lambda(k).
\]
Rearranging gives the \emph{excess-risk identity}
\begin{equation}\label{eq:B1-excess}
   \mathcal R^{\infty}(Q^{*})
   \;\le\;
   \mathcal R^{k}(Q) - \Lambda(k).
\end{equation}
Because the sequence is geometrically mixing,
Theorem 2 of \citet{mcdonald2011rademacher} (their bounded‐loss
Rademacher bound for $\alpha$–mixing processes) applies to
$\mathcal F_k$: with probability $\ge1-\delta$
\begin{equation}\label{eq:B1-gen}
   \mathcal R^{k}(Q)
   \;\le\;
   \hat{\mathcal R}^{k}(Q)
   +2\,\widehat{\Re}_n(\mathcal F_k)
   +3\,\frac{\ln(1/\delta)}{n}
   \quad\text{for \emph{every} }Q\in\mathcal H_k.
\end{equation}
Insert \eqref{eq:B1-gen} into the right–hand side of
\eqref{eq:B1-excess} to obtain \eqref{eq:B1-main}.\\
\end{proof}

\begin{corollary}[Oracle estimator and minimal order]
\label{cor:B2}
Assume the high-probability event of Proposition\;\ref{prop:B1}.
For a collection of pre-trained models
$\{Q_k\in\mathcal H_k\}_{k=1}^{M}$ define
\[
   \hat{\mathcal R}^{\infty}(Q^{*})
   :=\min_{1\le k\le M}
     \bigl\{\hat{\mathcal R}^{k}(Q_k)-\Lambda(k)\bigr\}.
\]
Then
\[
   \mathcal R^{\infty}(Q^{*})
   \;\le\;
   \hat{\mathcal R}^{\infty}(Q^{*})
   \quad\text{and}\quad
   k^{\dagger}\in\arg\min_{1\le k\le M}
      \bigl\{\hat{\mathcal R}^{k}(Q_k)-\Lambda(k)\bigr\}
   \Longrightarrow 
   k^{\dagger}\;\ge\;k^{*},
\]
where $k^{*}:=\inf\{k:\Lambda(k)=0\}$ is the minimal order for which the
universal learning curve vanishes.
\end{corollary}

\begin{proof}
Under the high-probability event,
inequality~\eqref{eq:B1-main} is valid for each
$k=1,\dots,M$ when evaluated at $Q_k$.
Discarding the non-negative complexity terms yields
\[
   \mathcal R^{\infty}(Q^{*})
   \;\le\;
   \hat{\mathcal R}^{k}(Q_k)-\Lambda(k),
   \qquad \forall k\le M.
\]
Taking the minimum over $k$ proves the claimed upper bound on
$\mathcal R^{\infty}(Q^{*})$.
If $k^{\dagger}$ realises that minimum while $k^{\dagger}<k^{*}$,
then $\Lambda(k^{\dagger})>\Lambda(k^{*})=0$ and monotonicity of
$\Lambda$ would contradict optimality of $k^{\dagger}$.  Hence
$k^{\dagger}\ge k^{*}$.
\end{proof}


\section{Estimation of $\mathbf{I}_{\text{pred}}$}
\label{appendix:partB}
\subsection{Estimating $\mathbf{I}_{\text{pred}}$ with Variational Neural Estimators}

We estimate the predictive mutual information $\mathbf{I}_{\text{pred}}$ using variational techniques based on recent advances in neural MI estimation. Given a $d$-dimensional time series $\{X_t\}$, our goal is to estimate
\[
\mathbf{I}_{\text{pred}}(\mathbf{X}_{t-k+1}^{t}; \mathbf{X}_{t+1}^{t+k'}),
\]
which quantifies the mutual information between a past window of $k$ time steps and a future window of $k'$ steps.

\paragraph{Variational Estimators.}
Our approach builds on contrastive lower bounds of mutual information, including SMILE~\cite{song2019understanding}, NWJ~\cite{nguyen2010estimating}, InfoNCE~\cite{oord2018}, TUBA~\cite{poole2019variational}, and DV~\cite{belghazi2018mine}. These methods rely on a parameterized critic function $f_\theta(x, y)$, which scores pairs of past and future segments. The critic is trained to distinguish between positive pairs (sampled from the joint distribution $p(x, y)$) and negative pairs (from the product of marginals $p(x)p(y)$).

\paragraph{Training Objective.}
The mutual information estimate is obtained by maximizing a variational objective of the form:
\begin{equation}
\mathbf{I}_{\text{pred}}(k, k') = \max_{\theta} \left\{
\frac{1}{B} \sum_{i=1}^{B} f_\theta(x_i, y_i)
- \ln \left( \frac{1}{B(B - 1)} \sum_{\substack{i,j=1 \\ i \ne j}}^B \exp(f_\theta(x_i, y_j)) \right)
\right\},
\label{eq:ipred}
\end{equation}
where $B$ is the batch size and the second term acts as a contrastive regularizer, penalizing high scores on mismatched pairs.

\paragraph{Critic Architectures and Optimization.}
To capture the structure of temporal data, we experiment with multiple critic architectures: \textit{Separable} (independent encodings for past and future), \textit{Concatenated} (joint embeddings), and \textit{Sequential} (LSTM-based encoders). The critic parameters $\theta$ are optimized using the Adam optimizer with stochastic gradient updates.

\paragraph{Data Sampling.}
Training batches are constructed either by sampling synthetic data from known generators, or by extracting context–future pairs from long, continuous sequences. The full training procedure is detailed in Algorithm~\ref{alg:evopred}.

\subsection{Synthetic data: kernel-based methods}\label{sub_section:synthetic_data_kernel}

For this part, we don't assume the invariance of temporal translation and stationarity of the process, as our concern is to verify that the estimators are working correctly.

\begin{proposition}
\label{prop:theoretical}(Theoretical value of $\mathbf{I}_{\text{pred}}$)
In the particular case where we consider $\{X_{t}\}_{t-k+1}^{t+k^\prime}$ a Gaussian process of dimension $d$ with all dimensions independent and of the same distribution, we can compute explicitly $\mathbf{I}_{\text{pred}}$.

$$\mathbf{I}_{\text{pred}}(k,k^\prime)=\mathbf{I}(\mathbf{X}_{t-k+1}^{t},{X}_{t+1}^{t+k^\prime})= \frac{d}{2} \ln \left(\frac{\left|\Sigma_1^{(1)}\right|\left|\Sigma_2^{(1)}\right|}{\left|\Sigma^{(1)}\right|}\right).$$

Where $\Sigma_1$ represents the covariance matrix of $\mathbf{X}_{t-k+1}^{t}$, $\Sigma_2$ the covariance matrix of ${X}_{t+1}^{t+k^\prime}$ and $\Sigma$ the joint covariance matrix. The index $(1)$ means that we can just look at the first dimension of the Gaussian process.
\end{proposition}

\begin{proof}

For a Gaussian process with dimension $d$ and independent dimensions, the predictive information formula is:
\begin{equation}
I(X_{\text{past}}; X_{\text{future}}) = \sum_{j=1}^d \frac{1}{2} \ln \left(\frac{|\Sigma_1^{(j)}||\Sigma_2^{(j)}|}{|\Sigma^{(j)}|}\right)
\end{equation}

Using the differential entropy property of a Gaussian vector $X$ of dimension $n$ with covariance matrix $\Sigma$:
\begin{equation}
h(X) = \frac{1}{2} \ln((2\pi e)^n |\Sigma|)
\end{equation}

Due to the independence of dimensions, the mutual information decomposes as:
\begin{equation}
I(X_{\text{past}}; X_{\text{future}}) = \sum_{j=1}^d I(X_{\text{past}}^{(j)}; X_{\text{future}}^{(j)})
\end{equation}

For each dimension $j$, with $\Sigma_1^{(j)}$ representing the covariance matrix of $X_{\text{past}}^{(j)}$, $\Sigma_2^{(j)}$ of $X_{\text{future}}^{(j)}$, and $\Sigma^{(j)}$ the joint covariance matrix:
\begin{align}
I(X_{\text{past}}^{(j)}; X_{\text{future}}^{(j)}) &= h(X_{\text{past}}^{(j)}) + h(X_{\text{future}}^{(j)}) - h(X_{\text{past}}^{(j)}, X_{\text{future}}^{(j)})\\
&= \frac{1}{2}\ln((2\pi e)^{n_1}|\Sigma_1^{(j)}|) + \frac{1}{2}\ln((2\pi e)^{n_2}|\Sigma_2^{(j)}|) - \frac{1}{2}\ln((2\pi e)^{n_1+n_2}|\Sigma^{(j)}|)\\
&= \frac{1}{2}\ln\left(\frac{|\Sigma_1^{(j)}||\Sigma_2^{(j)}|}{|\Sigma^{(j)}|}\right) + \frac{1}{2}\ln\left(\frac{(2\pi e)^{n_1}(2\pi e)^{n_2}}{(2\pi e)^{n_1+n_2}}\right)\\
&= \frac{1}{2}\ln\left(\frac{|\Sigma_1^{(j)}||\Sigma_2^{(j)}|}{|\Sigma^{(j)}|}\right)
\end{align}

Summing over all dimensions yields the result.
\end{proof}


 In practice, to ensure that the covariance matrix was invertible, we stayed with a low temporal resolution, from $k=5, k^\prime=10$ to $k=30, k^\prime=40$. If the covariance matrix was not invertible, we were just removing it from our results. To perform the experiences, we have chosen different Gaussian kernels, 
 \begin{itemize}
    \item \textbf{AR Kernel (Auto-Regressive):} $K_{\mathrm{AR}}(t_1,t_2) = \sigma^2 \, \rho^{\,|t_1-t_2|}$
    \item \textbf{Matérn 3/2 Kernel:} $ K_{\mathrm{M32}}(t_1,t_2) = \sigma^2 \left(1 + \frac{\sqrt{3}\,|t_1-t_2|}{l}\right) \exp\!\left(-\frac{\sqrt{3}\,|t_1-t_2|}{l}\right)$
  
    \item \textbf{Matérn 5/2 Kernel:} $K_{\mathrm{M52}}(t_1,t_2) = \sigma^2 \left(1 + \frac{\sqrt{5}\,|t_1-t_2|}{l} + \frac{5\,|t_1-t_2|^2}{3\,l^2}\right) \exp\!\left(-\frac{\sqrt{5}\,|t_1-t_2|}{l}\right)$
    
    \item \textbf{Squared Exponential Kernel:} $ K_{\mathrm{SE}}(t_1,t_2) = \sigma^2 \exp\!\left(-\frac{|t_1-t_2|^2}{2l^2}\right)$
    \item \textbf{Periodic Kernel:} $ K_{\mathrm{per}}(t_1,t_2) = \sigma^2 \exp\!\left(-\frac{2\,\sin^2\!\left(\pi\,\frac{|t_1-t_2|}{p}\right)}{l^2}\right)$
    \item \textbf{Rational Quadratic Kernel:} $  K_{\mathrm{RQ}}(t_1,t_2) = \sigma^2 \left(1 + \frac{|t_1-t_2|^2}{2\theta\,l^2}\right)^{-\theta}$
    \item \textbf{Locally Periodic Kernel:} $ K_{\mathrm{LP}}(t_1,t_2) = \sigma^2 \exp\!\left(-\frac{2\,\sin^2\!\left(\pi\,\frac{|t_1-t_2|}{p}\right)}{l^2}\right) \exp\!\left(-\frac{|t_1-t_2|^2}{2d^2}\right)$
\end{itemize}
For our experiment, we evaluated several variational lower bounds for mutual information estimation : $\hat{\mathbf{I}}_{\text{pred}}$--\texttt{SMILE}
~\cite{song2019understanding}, $\hat{\mathbf{I}}_{\text{pred}}$--\texttt{NWJ}
~\cite{nguyen2010estimating}, $\hat{\mathbf{I}}_{\text{pred}}$--\texttt{InfoNCE}
~\cite{oord2018representation}, $\hat{\mathbf{I}}_{\text{pred}}$--\texttt{DV}
~\cite{belghazi2018mine}, and $\hat{\mathbf{I}}_{\text{pred}}$--\texttt{TUBA}
~\cite{poole2019variational}. 
We took as parameters the followings, where in practice, $t_1$ and $t_2$ takes values between $0$ and $k+k\prime-1$. As we choose $t=k-1$ as the date of observation of $\{X_{t}\}_{t-k+1}^{t+k^\prime}$. The table \ref{tab:kernels_params} describe the parameters we took for each of the kernels in the experiments.


\begin{table}[H]
\centering
\begin{tabular}{lcccccc}
\hline
\textbf{Method} & $\rho$ & $l$ & \textbf{Period} & \textbf{Decay} & $\theta$ & $\sigma$ \\
\hline
AR                   & 0.8   & --  & --   & --   & --  & 0.5 \\
Matérn 3/2           & --    & 2.0 & --   & --   & --  & 1.0 \\
Matérn 5/2           & --    & 2.0 & --   & --   & --  & 1.0 \\
Squared Exponential  & --    & 2.0 & --   & --   & --  & 1.0 \\
Periodic             & --    & 3.0 & 2.0  & --   & --  & 0.5 \\
Rational Quadratic   & --    & 2.0 & --   & --   & 1.0 & 1.0 \\
Locally Periodic     & --    & 1.0 & 4.0  & 10.0 & --  & 1.0 \\
\hline
\end{tabular}
\caption{Parameters we choose for each of the kernels.}
\label{tab:kernels_params}
\end{table}

\subsection{Additional Combinations of $(k, k')$ for Predictive Information Estimation}
\label{sub_section:additionnal_combinations}

This section presents complementary estimation results for the predictive information $\mathbf{I}_{\text{pred}}(k, k')$ obtained with various combinations of past and future context lengths $(k, k')$. These results extend the main analysis shown in Figure~\ref{fig:estimation_evoPred}, where $k=5$ and $k'=10$ were fixed. 

\textbf{Case $k = 2$, $k' = 5$:}
\begin{figure}[H]
    \centering
    \includegraphics[width=1\linewidth]{figures/k_2_k_prime_5.png} 
    \caption{
        Estimation of $\mathbf{I}_{\text{pred}}(2, 5)$ using various neural-based methods.
    }
    \label{fig:estimation_evoPred_k_2_k_prime_5}
\end{figure}

\textbf{Case $k = 10$, $k' = 20$:}
\begin{figure}[H]
    \centering
    \includegraphics[width=1\linewidth]{figures/k_10_k_prime_20.png} 
    \caption{
        Estimation of $\mathbf{I}_{\text{pred}}(10, 20)$ using various neural-based methods.
    }
    \label{fig:estimation_evoPred_k_10_k_prime_20}
\end{figure}
% Figure~\ref{fig:estimation_evoPred} illustrates the estimation results of $\mathbf{I}_{\text{pred}}(k, k^\prime)$ across various input dimensions and kernel types. We evaluated several variational lower bounds for mutual information estimation, including the Smoothed Mutual Information Lower-bound Estimator ($\hat{\mathbf{I}}_{\text{pred}}$--\texttt{SMILE}
% )~\cite{song2019understanding}, the Nguyen–Wainwright–Jordan bound ($\hat{\mathbf{I}}_{\text{pred}}$--\texttt{NWJ}
% )~\cite{nguyen2010estimating}, Contrastive Predictive Coding ($\hat{\mathbf{I}}_{\text{pred}}$--\texttt{InfoNCE}
% )~\cite{oord2018representation}, the Donsker–Varadhan bound ($\hat{\mathbf{I}}_{\text{pred}}$--\texttt{DV}
% )~\cite{belghazi2018mine}, and the Tractable Unnormalized Barber–Agakov bound ($\hat{\mathbf{I}}_{\text{pred}}$--\texttt{TUBA}
% )~\cite{poole2019variational}. $\mathbf{I}_{\text{pred}}$--\texttt{True} is the theoretical value calculated with ~\ref{prop:theoretical}.

% \begin{figure}[H]
%     \centering
%     \includegraphics[width=1\linewidth]{figures/output.png} 
%     \captionsetup{font=small}
%     \caption{Estimation of $\mathbf{I}_{\text{pred}}(k, k')$ using various neural-based methods. Blue regions indicate negative bias (underestimation), while red regions indicate positive bias (overestimation); the color intensity reflects the magnitude of the bias. For this experiment, we fixed $k=5$ and $k'=10$. All methods were implemented using a concatenated architecture, as described above. }
%     \label{fig:estimation_evoPred}
% \end{figure}

% While all methods provide relatively accurate estimates in low-dimensional settings (dimensionality $d \leq 20$), their performance deteriorates as the input space becomes more complex. Notably, methods like evoPred-SMILE and evoPred-NWJ exhibited consistent underestimation, in structured, high-dimensional regimes (e.g., Periodic and RQ kernels). Although our estimator of $\mathbf{I}_{\text{pred}}$ is not perfect, it is important to emphasize that the primary objective of this paper is not to achieve precise estimation of predictive information, but rather to propose a method for determining the optimal forecasting risk.


\subsection{Other simple process}\label{subsection:gaussian_process_toy}
Consider the process \( \{z_t\}_{t=0}^{N-1} \subset \mathbb{R}^d \) defined as follows:
\begin{align}
    z_0 &\sim \mathcal{N}(0, I_d), \\
    z_t &= \rho\, z_{t-1} + \sqrt{1 - \rho^2}\,\varepsilon_t,\quad t \geq 1,
\end{align}
where \( \rho \in (-1,1) \) is the correlation coefficient, \( I_d \) is the \( d \times d \) identity matrix, and \( \varepsilon_t \sim \mathcal{N}(0, I_d) \) are i.i.d. random vectors.


This defines a Markov (AR(1)) Gaussian process in which each state depends solely on the immediate predecessor. If the total time length is given by \( N = T_{\text{past}} + T_{\text{fut}} \), we partition the sequence into:
\begin{align*}
X_{\text{past}} &= \{ z_0, z_1, \dots, z_{T_{\text{past}} - 1} \}, \\
X_{\text{fut}} &= \{ z_{T_{\text{past}}}, \dots, z_{N-1} \}.
\end{align*}

Then we compute the theorical value of $\mathbf{I}_{\text{pred}}$ thanks to \ref{prop:markov_combined}.

\subsection{Estimator comparison for $T_{past} = 30$, $T_{future} = 40$ }

Below are the estimated values of $\mathbf{I}_{\text{pred}}$ across different kernel types, critic architectures, and $\rho$ values. The training parameters are as follows: a batch size of 70, 10{,}000 iterations, and a learning rate of $5 \times 10^{-4}$.


\begin{table}[H]
\centering
\resizebox{\textwidth}{!}{%
\begin{tabular}{cccccccc}
\toprule
dim & theoretical\_mi & rho & SeparableCritic & ConcatCritic & SequentialCritic & EvoRate & kernel\_type \\
\midrule
1   & 0.14  & 0.50 & 0.14  & 0.14  & 0.01  & \textbf{0.14} &  \\
1   & 0.51  & 0.80 & 0.48  & \textbf{0.51}  & 0.28  & 0.48  &  \\
1   & 0.83  & 0.90 & 0.77  & \textbf{0.82}  & 0.44  & 0.79  &  \\
5   & 0.72  & 0.50 & 0.63  & \textbf{0.67}  & -0.07 & 0.65  &  \\
5   & 2.55  & 0.80 & 2.16  & \textbf{2.41}  & 1.37  & 2.13  &  \\
5   & 4.15  & 0.90 & 3.17  & \textbf{3.75}  & 2.15  & 3.34  &  \\
20  & 2.88  & 0.50 & 0.60  & \textbf{1.45}  & -1.79 & 0.60  &  \\
20  & 10.22 & 0.80 & 3.21  & \textbf{6.07}  & 5.86  & 3.16  &  \\
20  & 16.61 & 0.90 & 4.97  & \textbf{9.28}  & 8.87  & 4.95  &  \\
100 & 14.38 & 0.50 & 0.41  & 0.78  & \textbf{5.14}  & 0.51  &  \\
100 & 51.08 & 0.80 & 2.13  & 3.76  & \textbf{45.64} & 2.46  &  \\
100 & 83.04 & 0.90 & 4.40  & 5.46  & \textbf{63.96} & 4.58  &  \\
1   & 0.51  &      & 0.48  & \textbf{0.50}  & 0.28  & 0.47  & AR \\
1   & 0.62  &      & 0.58  & \textbf{0.59}  & 0.23  & 0.57  & matern32 \\
1   &       &      & -6.83 & 7.09  & 0.58  & 1.39  & periodic \\
5   & 2.55  &      & 1.83  & \textbf{2.39}  & 1.10  & 2.22  & AR \\
5   & 3.09  &      & 2.52  & \textbf{2.75}  & 0.97  & 2.62  & matern32 \\
5   &       &      & 7.12  & 28.55 & 2.83  & 3.93  & periodic \\
20  & 10.22 &      & 2.12  & \textbf{4.92}  & 4.08  & 2.11  & AR \\
20  & 12.35 &      & 2.33  & 5.40  & \textbf{5.45}  & 3.23  & matern32 \\
20  &       &      & 13.34 & 31.06 & 11.61 & 4.96  & periodic \\
100 & 51.08 &      & 2.40  & 3.63  & \textbf{20.77} & 1.87  & AR \\
100 & 61.74 &      & 2.07  & 3.06  & \textbf{44.52} & 2.37  & matern32 \\
100 &       &      & 9.47  & 26.18 & 81.52 & 7.40  & periodic \\
\bottomrule
\end{tabular}%
}
\caption{Comparison of predictive information estimators across different kernel types and process dimensions. Values rounded to two decimals.}
\label{tab:mesures}
\end{table}


% \subsection{Estimator comparison $T_{past}=5$ and $T_{future}=10$}

% This configuration sets up an experiment to run on CPU, using two data configurations: one employing a Markov Gaussian model with dimensions $[1, 5, 20, 100]$ and rho values $[0.5, 0.8, 0.9]$, and another employing a kernel-based model with the same dimensions and specific kernel parameters for AR (rho: 0.8, sigma: 0.5), matern32 (l: 2.0, sigma: 1.0), and periodic (alpha: 0.2, sigma: 1.0); it also defines several estimators, enabling SeparableCriticSMILE, ConcatCriticSMILE, SequentialCriticSMILE, and EvoRateSeparable while disabling SeparableCriticNWJ and SeparableCriticInfoNCE, and specifies time windows ($T_{past}=5$ and $T_{future}=10$) alongside training settings (a batch size of 70, 10,000 iterations, and a learning rate of 5e-4), with results being saved in the "results" directory.
% \begin{table}[H]
% \centering
% \footnotesize
% \begin{tabular}{cccccccc}
% \toprule
% dim & theoreticalMI & rho & SeparableCritic & ConcatCritic & SequentialCritic & EvoRate & kernel\_type \\
% \midrule
% 1   & 0.14  & 0.50 & 0.14  & 0.14  & -0.01 & \textbf{0.14} & \\
% 1   & 0.51  & 0.80 & 0.48  & \textbf{0.52}  & 0.30  & 0.50  & \\
% 1   & 0.83  & 0.90 & 0.79  & \textbf{0.81}  & 0.41  & 0.78  & \\
% 5   & 0.72  & 0.50 & 0.67  & \textbf{0.69}  & -0.57 & 0.65  & \\
% 5   & 2.55  & 0.80 & 2.27  & \textbf{2.52}  & 1.42  & 2.31  & \\
% 5   & 4.15  & 0.90 & 3.42  & \textbf{3.87}  & 2.13  & 3.58  & \\
% 20  & 2.88  & 0.50 & 1.37  & \textbf{1.58}  & -1.52 & 1.36  & \\
% 20  & 10.22 & 0.80 & 4.87  & \textbf{6.68}  & 5.87  & 4.84  & \\
% 20  & 16.61 & 0.90 & 8.18  & \textbf{10.53} & 8.72  & 7.48  & \\
% 100 & 14.38 & 0.50 & 1.05  & 1.73  & \textbf{5.12}  & 1.03  & \\
% 100 & 51.08 & 0.80 & 3.99  & 5.31  & \textbf{46.86} & 4.35  & \\
% 100 & 83.04 & 0.90 & 6.40  & 8.01  & \textbf{64.62} & 6.33  & \\
% 1   & 0.51  &       & 0.49  & \textbf{0.50}  & 0.28  & 0.50  & AR \\
% 1   & 0.62  &       & 0.58  & \textbf{0.59}  & 0.24  & 0.58  & matern32 \\
% 1   & 7.51  &       & 0.93  & \textbf{5.19}  & 0.62  & 2.64  & periodic \\
% 5   & 2.55  &       & 2.25  & \textbf{2.47}  & 1.14  & 2.25  & AR \\
% 5   & 3.09  &       & 2.69  & \textbf{2.88}  & 1.37  & 2.69  & matern32 \\
% 5   & 37.55 &       & 2.23  & \textbf{31.47} & 2.85  & 8.46  & periodic \\
% 20  & 10.22 &       & 3.41  & \textbf{5.66}  & 4.09  & 3.37  & AR \\
% 20  & 12.35 &       & 4.96  & \textbf{5.93}  & 5.65  & 4.96  & matern32 \\
% 20  & 150.19&       & 11.86 & \textbf{31.32} & 11.50 & 12.40 & periodic \\
% 100 & 51.08 &       & 3.22  & 5.22  & \textbf{20.82} & 2.89  & AR \\
% 100 & 61.74 &       & 3.85  & 5.77  & \textbf{44.97} & 3.68  & matern32 \\
% 100 & 750.97&       & 5.76  & 25.85 & \textbf{81.14} & 8.57  & periodic \\
% \bottomrule
% \end{tabular}
% \caption{Results with values rounded to two decimals; the estimator closest to the theoretical\_mi in each row is highlighted in bold.}
% \label{tab:results}
% \end{table}

\section{Experiment}
\label{sec:experiment}

\subsubsection*{Note on the Use of Nats}
\label{subsec:nats_explanation}

\noindent
\textbf{Units in nats.} Throughout our experiments, all information-theoretic results and measures (e.g., mutual information and differential entropy) are reported in \emph{nats} rather than in bits. In information theory, the choice of base for the logarithm determines the unit: base-$e$ (the natural logarithm) yields nats, while base-$2$ yields bits. We prefer natural logarithms because they often simplify analytical expressions in both theory and implementation (e.g., when computing the log-likelihood in many machine learning frameworks). However, it is straightforward to convert from nats to bits by noting
\[
1 \, \text{nat} \;=\; \frac{1}{\ln(2)} \;\text{bits}.
\]
Due to this simple relationship, one can easily switch to bits by scaling the values by $1 / \ln(2)$ if desired.

\textbf{Estimation of $\tilde{\Lambda}(k)$}
\label{sub_section:learning_curve}

To estimate the learning curve $\tilde{\Lambda}(k) = l(k) - l_0$, we evaluate the conditional entropy rate $l(k)$ from the data and analytically compute the theoretical baseline $l_0$. From the probabilistic formulation in Equation~\eqref{eq:AR6}, we can directly access the data distribution, which allows us to compute the entropy as follows:

\paragraph{Estimation of $l(k)$.}  

We estimate the conditional entropy \(l(k) = H(X_t \mid X_{t-k}, \ldots, X_{t-1})\) of a multivariate time series using ridge regression. After fitting a linear model to predict \(X_t\) from its past, we compute the residuals and estimate their empirical covariance matrix \(\Sigma\). Assuming the residuals are approximately Gaussian, the conditional entropy is estimated using:
\[
l(k) \approx \frac{1}{2} \ln\left( (2\pi e)^d \cdot |\Sigma| \right),
\]
where \(d\) is the dimension of the observed vectors. 

\textbf{Estimation of $l_0$.}  

According to Proposition~\ref{prop:markov_combined}, we have $l_0 = l(p)$ for an autoregressive process of order $p$. Furthermore, from Equation~\ref{eq:AR6},  $X_t \mid X_{t-1}, \ldots, X_{t-p}$ follows a normal distribution with 
\[
\mu = \frac{\rho}{p} \sum_{j=t-p}^{t-1} X_j
\quad \text{and} \quad
\sigma^2 = 1 - \rho^2.
\]
Since $l(p) = H(X_t \mid X_{t-1}, \ldots, X_{t-p})$, we can compute it directly using the above formula for the conditional entropy of a multivariate normal distribution.

% \begin{table}[H]
%     \centering
%     \begin{tabular}{l c c c c}
%         \toprule
%         Model & \(k\) & Final Loss (mean last 100) & Universal Learning Rate \(\lambda(k)\) & Loss - \(\lambda(k)\) \\
%         \midrule
%         \multicolumn{5}{c}{\textbf{LSTM}} \\
%         \midrule
%         LSTM & 1  & 0.9824 & 0.3887 & 0.5937 \\
%         LSTM & 2  & 0.8946 & 0.2663 & 0.6283 \\
%         LSTM & 3  & 0.8267 & 0.2074 & 0.6193 \\
%         LSTM & 4  & 0.8584 & 0.1726 & 0.6858 \\
%         LSTM & 5  & 0.7554 & 0.1491 & 0.6063 \\
%         \midrule
%         LSTM & 15 & 0.6523 & 0.0729 & 0.5794 \\
%         LSTM & 16 & 0.7939 & 0.0680 & 0.7258 \\
%         LSTM & 17 & 0.6624 & 0.0610 & 0.6014 \\
%         LSTM & 18 & 0.7128 & 0.0489 & 0.6639 \\
%         LSTM & 19 & 0.7600 & 0.0240 & 0.7360 \\
%         \midrule
%         \multicolumn{5}{c}{\textbf{MLP}} \\
%         \midrule
%         MLP  & 1  & 1.0003 & 0.3887 & 0.6116 \\
%         MLP  & 2  & 0.8691 & 0.2663 & 0.6028 \\
%         MLP  & 3  & 0.7705 & 0.2074 & 0.5630 \\
%         MLP  & 4  & 0.7421 & 0.1726 & 0.5695 \\
%         MLP  & 5  & 0.7511 & 0.1491 & 0.6020 \\
%         \midrule
%         MLP  & 15 & 0.7433 & 0.0729 & 0.6704 \\
%         MLP  & 16 & 0.7296 & 0.0680 & 0.6616 \\
%         MLP  & 17 & 0.6890 & 0.0610 & 0.6279 \\
%         MLP  & 18 & 0.7216 & 0.0489 & 0.6727 \\
%         MLP  & 19 & 0.6806 & 0.0240 & 0.6566 \\
%         \midrule
%         \multicolumn{5}{c}{\textbf{Transformer}} \\
%         \midrule
%         Transformer & 1  & 0.9636 & 0.3887 & 0.5749 \\
%         Transformer & 2  & 0.8724 & 0.2663 & 0.6061 \\
%         Transformer & 3  & 0.8159 & 0.2074 & 0.6085 \\
%         Transformer & 4  & 0.7224 & 0.1726 & 0.5499 \\
%         Transformer & 5  & 0.7840 & 0.1491 & 0.6350 \\
%         \midrule
%         Transformer & 15 & 0.7462 & 0.0729 & 0.6733 \\
%         Transformer & 16 & 0.7468 & 0.0680 & 0.6787 \\
%         Transformer & 17 & 0.7485 & 0.0610 & 0.6875 \\
%         Transformer & 18 & 0.7120 & 0.0489 & 0.6630 \\
%         Transformer & 19 & 0.6453 & 0.0240 & 0.6213 \\
%         \bottomrule
%     \end{tabular}
%     \caption{Comparison of Final Loss (mean last 100), Universal Learning Rate \(\lambda(k)\), and Loss - \(\lambda(k)\) for different regression orders \(k\) (showing \(k\in\{1,2,3,4,5\}\) and \(k\in\{15,\ldots,19\}\)) for the LSTM, MLP, and Transformer models.}
%     \label{tab:loss_comparison_subset}
% \end{table}

\subsection{Ising Spin Sequences}
\label{sec:partC}

\textbf{Training of the MLP and LSTM models.}
We train both a multilayer perceptron (MLP) and a long short-term memory (LSTM) model on spin chain data generated using the procedure described in Section~\ref{sec:ising}. Each model receives a context window of length $k$ and is trained to predict the next binary symbol in the sequence.

The MLP model consists of two hidden layers with ReLU activations, mapping the input vector of length $k$ to a softmax output over two classes. The LSTM model, on the other hand, processes the input sequence as a series of scalar values through an LSTM layer followed by a fully connected output layer.

Training is performed using the Adam optimizer with a learning rate of $10^{-3}$ and a batch size of 128. For each value of $k$, training proceeds for up to 1000 epochs, using early stopping with a patience of 10 epochs based on validation loss. To standardize comparisons across models, we fix the number of batches per epoch and apply the same evaluation procedure to all architectures. The dataset is split into 80\% training and 20\% validation sets, and results are averaged across runs to account for variance.

{\textbf{Model architectures:}}
\begin{itemize}
    \item {\textbf{MLP}: 2 hidden layers (64, then 32 neurons), ReLU activations, output layer with 2 units (softmax).}
    \item {\textbf{LSTM}: Single-layer LSTM (32 hidden units), fully connected layer mapping to 2 output classes.}
\end{itemize}

{\textbf{Training setup:} Adam optimizer (lr=$10^{-3}$), batch size=128, max epochs=1000, early stopping (patience=10), 50 batches per epoch, 80/20 train-validation split.}\\

Below are the results obtained for different block sizes \textit{M}: (1) 10{,}000, (2) 100{,}000, (3) 1{,}000{,}000, and (4) 10{,}000{,}000.

\begin{table}[H]
    \centering
     \resizebox{\textwidth}{!}{%
    \begin{tabular}{c c c c c c c}
        \toprule
        $k$ & $\hat{\mathcal{R}}^{k}_{\text{LSTM}}(Q)$ & $\hat{\mathcal{R}}^{k}_{\text{MLP}}(Q)$ & $\hat{\Lambda}(k)$ & 
        $\hat{\mathcal{R}}^{k}_{\text{LSTM}}(Q) - \hat{\Lambda}(k)$ & 
        $\hat{\mathcal{R}}^{k}_{\text{MLP}}(Q) - \hat{\Lambda}(k)$ & 
        \texttt{EvoRate}$(k)$ \\
        \midrule
        1  & 0.6932 & 0.6927 & 0.3208 ± 0.0043 & \textbf{0.3724 ± 0.0043} & \textbf{0.3719 ± 0.0043} & 0.0005 ± 0.0012 \\
        2  & 0.6220 & 0.6200 & 0.2058 ± 0.0023 & 0.4162 ± 0.0023 & 0.4142 ± 0.0023 & 0.1171 ± 0.0080 \\
        3  & 0.5955 & 0.5571 & 0.1503 ± 0.0031 & 0.4452 ± 0.0031 & 0.4068 ± 0.0031 & 0.1731 ± 0.0105 \\
        4  & 0.5611 & 0.5560 & 0.1177 ± 0.0034 & 0.4434 ± 0.0034 & 0.4383 ± 0.0034 & 0.2059 ± 0.0115 \\
        5  & 0.5394 & 0.5429 & 0.0962 ± 0.0035 & 0.4432 ± 0.0035 & 0.4467 ± 0.0035 & 0.2274 ± 0.0120 \\
        6  & 0.5347 & 0.5343 & 0.0810 ± 0.0036 & 0.4537 ± 0.0036 & 0.4533 ± 0.0036 & 0.2427 ± 0.0122 \\
        7  & 0.5167 & 0.5199 & 0.0696 ± 0.0035 & 0.4471 ± 0.0035 & 0.4503 ± 0.0035 & 0.2541 ± 0.0124 \\
        8  & 0.5105 & 0.5088 & 0.0607 ± 0.0035 & 0.4498 ± 0.0035 & 0.4481 ± 0.0035 & 0.2629 ± 0.0125 \\
        9  & 0.5036 & 0.5092 & 0.0537 ± 0.0034 & 0.4499 ± 0.0034 & 0.4555 ± 0.0034 & 0.2700 ± 0.0126 \\
        10 & 0.5079 & 0.5064 & 0.0479 ± 0.0034 & 0.4600 ± 0.0034 & 0.4585 ± 0.0034 & 0.2758 ± 0.0126 \\
        11 & 0.5077 & 0.5232 & 0.0430 ± 0.0034 & 0.4647 ± 0.0034 & 0.4802 ± 0.0034 & 0.2807 ± 0.0127 \\
        12 & 0.5181 & 0.4997 & 0.0389 ± 0.0034 & 0.4792 ± 0.0034 & 0.4608 ± 0.0034 & 0.2850 ± 0.0127 \\
        13 & 0.4965 & 0.5012 & 0.0354 ± 0.0033 & 0.4611 ± 0.0033 & 0.4658 ± 0.0033 & 0.2887 ± 0.0128 \\
        14 & 0.4989 & 0.4992 & 0.0324 ± 0.0033 & 0.4665 ± 0.0033 & 0.4668 ± 0.0033 & 0.2924 ± 0.0128 \\
        15 & 0.4864 & 0.5000 & 0.0297 ± 0.0032 & 0.4567 ± 0.0032 & 0.4703 ± 0.0032 & 0.2962 ± 0.0128 \\
        16 & 0.4949 & 0.5062 & 0.0273 ± 0.0032 & 0.4676 ± 0.0032 & 0.4789 ± 0.0032 & 0.3010 ± 0.0128 \\
        17 & 0.4972 & 0.4968 & 0.0252 ± 0.0032 & 0.4720 ± 0.0032 & 0.4716 ± 0.0032 & 0.3080 ± 0.0128 \\
        18 & 0.4885 & \textbf{0.4872} & 0.0232 ± 0.0033 & 0.4653 ± 0.0033 & 0.4640 ± 0.0033 & 0.3202 ± 0.0128 \\
        19 & \textbf{0.4853} & 0.5041 & 0.0213 ± 0.0034 & 0.4640 ± 0.0034 & 0.4828 ± 0.0034 & 0.3452 ± 0.0129 \\
        \midrule
        \textbf{Minima} & \textbf{0.4853} & \textbf{0.4872} & -- & $\hat{\mathcal{R}}^\infty(Q^*) = 0.3724 \pm 0.0043$ & $\hat{\mathcal{R}}^\infty(Q^*) = 0.3719 \pm 0.0043$ & -- \\
        \bottomrule
    \end{tabular}
    }
    \caption{Estimated risks and $\Lambda(k)$ values for different history lengths $k$ (\textit{M} = 10{,}000).}
    \label{tab:universal_learning_1}
\end{table}

\begin{table}[H]
\centering
 \resizebox{\textwidth}{!}{%
\begin{tabular}{c c c c c c c}
\toprule
Order $k$ & $\hat{\mathcal{R}}^{k}_{\text{LSTM}}(Q)$ & $\hat{\mathcal{R}}^{k}_{\text{MLP}}(Q)$ & $\hat{\Lambda}(k)$ & 
$\hat{\mathcal{R}}^{k}_{\text{LSTM}}(Q) - \hat{\Lambda}(k)$ & $\hat{\mathcal{R}}^{k}_{\text{MLP}}(Q) - \hat{\Lambda}(k)$ & \texttt{EvoRate}$(k)$ \\
\midrule
1  & 0.6897 & 0.6933 & 0.3213 ± 0.0153 & \textbf{0.3684 ± 0.0153} & \textbf{0.3720 ± 0.0153} & 0.0056 ± 0.0030 \\
2  & 0.5789 & 0.6070 & 0.2061 ± 0.0058 & 0.3728 ± 0.0058 & 0.4009 ± 0.0058 & 0.1261 ± 0.0208 \\
3  & 0.6036 & 0.5745 & 0.1506 ± 0.0071 & 0.4530 ± 0.0071 & 0.4239 ± 0.0071 & 0.1832 ± 0.0262 \\
4  & 0.5497 & 0.5615 & 0.1179 ± 0.0081 & 0.4318 ± 0.0081 & 0.4436 ± 0.0081 & 0.2163 ± 0.0282 \\
5  & 0.5146 & 0.5383 & 0.0964 ± 0.0085 & 0.4182 ± 0.0085 & 0.4419 ± 0.0085 & 0.2379 ± 0.0290 \\
6  & 0.5226 & 0.5537 & 0.0811 ± 0.0087 & 0.4415 ± 0.0087 & 0.4726 ± 0.0087 & 0.2531 ± 0.0294 \\
7  & 0.5311 & 0.5409 & 0.0697 ± 0.0088 & 0.4614 ± 0.0088 & 0.4712 ± 0.0088 & 0.2645 ± 0.0296 \\
8  & 0.5124 & 0.5531 & 0.0608 ± 0.0088 & 0.4516 ± 0.0088 & 0.4923 ± 0.0088 & 0.2733 ± 0.0296 \\
9  & 0.5236 & 0.5331 & 0.0537 ± 0.0088 & 0.4699 ± 0.0088 & 0.4794 ± 0.0088 & 0.2803 ± 0.0297 \\
10 & 0.4901 & 0.5186 & 0.0479 ± 0.0088 & 0.4422 ± 0.0088 & 0.4707 ± 0.0088 & 0.2861 ± 0.0297 \\
11 & \textbf{0.4679} & 0.5119 & 0.0430 ± 0.0088 & 0.4249 ± 0.0088 & 0.4689 ± 0.0088 & 0.2910 ± 0.0296 \\
12 & 0.4975 & 0.5273 & 0.0389 ± 0.0088 & 0.4586 ± 0.0088 & 0.4884 ± 0.0088 & 0.2952 ± 0.0296 \\
13 & 0.4692 & 0.4736 & 0.0353 ± 0.0088 & 0.4339 ± 0.0088 & 0.4383 ± 0.0088 & 0.2990 ± 0.0296 \\
14 & 0.5152 & 0.5259 & 0.0323 ± 0.0088 & 0.4829 ± 0.0088 & 0.4936 ± 0.0088 & 0.3026 ± 0.0296 \\
15 & 0.4971 & 0.4733 & 0.0296 ± 0.0087 & 0.4675 ± 0.0087 & 0.4437 ± 0.0087 & 0.3064 ± 0.0296 \\
16 & 0.4714 & 0.5184 & 0.0272 ± 0.0087 & 0.4442 ± 0.0087 & 0.4912 ± 0.0087 & 0.3112 ± 0.0296 \\
17 & 0.4996 & 0.5278 & 0.0250 ± 0.0087 & 0.4746 ± 0.0087 & 0.5028 ± 0.0087 & 0.3182 ± 0.0295 \\
18 & 0.5037 & 0.4748 & 0.0230 ± 0.0089 & 0.4807 ± 0.0089 & 0.4518 ± 0.0089 & 0.3305 ± 0.0297 \\
19 & 0.5189 & \textbf{0.4668} & 0.0211 ± 0.0092 & 0.4978 ± 0.0092 & 0.4457 ± 0.0092 & 0.3556 ± 0.0299 \\
\midrule
\multicolumn{1}{r}{\textbf{Minima}} &
\textbf{0.4679} &
\textbf{0.4668} &
-- &
$\hat{\mathcal{R}}^\infty(Q^*) = 0.3684 \pm 0.0153$ &
$\hat{\mathcal{R}}^\infty(Q^*) = 0.3720 \pm 0.0153$ &
-- \\
\bottomrule
\end{tabular}
}
\caption{Estimated risks and $\Lambda(k)$ values for different history lengths $k$ (\textit{M} = 100{,}000).}
\label{tab:universal_learning_2}
\end{table}


\begin{table}[H]
\centering
 \resizebox{\textwidth}{!}{%
\begin{tabular}{c c c c c c c}
\toprule
Order $k$ & $\hat{\mathcal{R}}^{k}_{\text{LSTM}}(Q)$ & $\hat{\mathcal{R}}^{k}_{\text{MLP}}(Q)$ & $\hat{\Lambda}(k)$ & 
$\hat{\mathcal{R}}^{k}_{\text{LSTM}}(Q) - \hat{\Lambda}(k)$ & $\hat{\mathcal{R}}^{k}_{\text{MLP}}(Q) - \hat{\Lambda}(k)$ & \texttt{EvoRate}$(k)$ \\
\midrule
1  & 0.6910 & 0.6808 & 0.3235 ± 0.0718 & 0.3675 ± 0.0718 & 0.3573 ± 0.0718 & 0.0009 ± 0.0320 \\
2  & 0.5641 & 0.6290 & 0.2069 ± 0.0188 & 0.3572 ± 0.0188 & 0.4221 ± 0.0188 & 0.1450 ± 0.0465 \\
3  & 0.6157 & 0.5502 & 0.1510 ± 0.0314 & 0.4647 ± 0.0314 & 0.3992 ± 0.0314 & 0.2090 ± 0.0554 \\
4  & 0.5913 & 0.5847 & 0.1181 ± 0.0376 & 0.4732 ± 0.0376 & 0.4666 ± 0.0376 & 0.2460 ± 0.0588 \\
5  & 0.5621 & 0.5145 & 0.0965 ± 0.0395 & 0.4656 ± 0.0395 & 0.4180 ± 0.0395 & 0.2706 ± 0.0603 \\
6  & 0.5087 & 0.6114 & 0.0811 ± 0.0402 & 0.4276 ± 0.0402 & 0.5303 ± 0.0402 & 0.2883 ± 0.0609 \\
7  & 0.4705 & 0.4423 & 0.0697 ± 0.0401 & 0.4008 ± 0.0401 & 0.3726 ± 0.0401 & 0.3016 ± 0.0613 \\
8  & 0.5378 & 0.5469 & 0.0607 ± 0.0398 & 0.4771 ± 0.0398 & 0.4862 ± 0.0398 & 0.3120 ± 0.0614 \\
9  & 0.5195 & 0.5008 & 0.0536 ± 0.0394 & 0.4659 ± 0.0394 & 0.4472 ± 0.0394 & 0.3202 ± 0.0615 \\
10 & 0.4480 & 0.4690 & 0.0478 ± 0.0390 & 0.4002 ± 0.0390 & 0.4212 ± 0.0390 & 0.3269 ± 0.0615 \\
11 & 0.5612 & 0.5034 & 0.0429 ± 0.0386 & 0.5183 ± 0.0386 & 0.4605 ± 0.0386 & 0.3326 ± 0.0615 \\
12 & 0.5172 & 0.4650 & 0.0388 ± 0.0383 & 0.4784 ± 0.0383 & 0.4262 ± 0.0383 & 0.3373 ± 0.0615 \\
13 & 0.5039 & 0.5169 & 0.0352 ± 0.0379 & 0.4687 ± 0.0379 & 0.4817 ± 0.0379 & 0.3416 ± 0.0615 \\
14 & 0.5588 & 0.4728 & 0.0321 ± 0.0376 & 0.5267 ± 0.0376 & 0.4407 ± 0.0376 & 0.3456 ± 0.0615 \\
15 & 0.4760 & 0.4402 & 0.0294 ± 0.0374 & 0.4466 ± 0.0374 & 0.4108 ± 0.0374 & 0.3498 ± 0.0614 \\
16 & 0.5425 & \textbf{0.3660} & 0.0270 ± 0.0373 & 0.5155 ± 0.0373 & \textbf{0.3390 ± 0.0373} & 0.3548 ± 0.0614 \\
17 & 0.4003 & 0.4738 & 0.0249 ± 0.0370 & 0.3754 ± 0.0370 & 0.4489 ± 0.0370 & 0.3619 ± 0.0616 \\
18 & \textbf{0.3798} & 0.4653 & 0.0229 ± 0.0359 & \textbf{0.3569 ± 0.0359} & 0.4424 ± 0.0359 & 0.3741 ± 0.0618 \\
19 & 0.5082 & 0.5629 & 0.0210 ± 0.0315 & 0.4872 ± 0.0315 & 0.5419 ± 0.0315 & 0.3989 ± 0.0618 \\
\midrule
\multicolumn{1}{r}{\textbf{Minima}} &
\textbf{0.3798} &
\textbf{0.3660} &
-- &
$\hat{\mathcal{R}}^\infty(Q^*) = 0.3569 \pm 0.0359$ &
$\hat{\mathcal{R}}^\infty(Q^*) = 0.3390 \pm 0.0373$ &
-- \\
\bottomrule
\end{tabular}
}
\captionsetup{justification=centering}
\caption{Estimated risks and $\Lambda(k)$ values for different history lengths $k$ (\textit{M} = 1{,}000{,}000).}
\label{tab:universal_learning_3}
\end{table}

\begin{table}[H]
\centering
\resizebox{\textwidth}{!}{%
\begin{tabular}{c c c c c c c}
\toprule
Order $k$ & $\hat{\mathcal{R}}^{k}_{\text{LSTM}}(Q)$ & $\hat{\mathcal{R}}^{k}_{\text{MLP}}(Q)$ & $\hat{\Lambda}(k)$ & 
$\hat{\mathcal{R}}^{k}_{\text{LSTM}}(Q) - \hat{\Lambda}(k)$ & $\hat{\mathcal{R}}^{k}_{\text{MLP}}(Q) - \hat{\Lambda}(k)$ & \texttt{EvoRate}(k) \\
\midrule
1  & 0.6709 & 0.3551 & 0.0037 ± 0.0929 & 0.6672 ± 0.0929 & 0.3514 ± 0.0929 & 0.4760 ± 0.2277 \\
2  & 0.5209 & 0.2713 & 0.0037 ± 0.0067 & 0.5172 ± 0.0067 & 0.2676 ± 0.0067 & 0.4760 ± 0.2277 \\
3  & 0.1961 & 0.6698 & 0.0037 ± 0.0086 & 0.1924 ± 0.0086 & 0.6661 ± 0.0086 & 0.4760 ± 0.2277 \\
4  & 0.5440 & 0.6037 & 0.0037 ± 0.0075 & 0.5403 ± 0.0075 & 0.6000 ± 0.0075 & 0.4760 ± 0.2277 \\
5  & 0.2271 & 0.1029 & 0.0037 ± 0.0108 & 0.2234 ± 0.0108 & 0.0992 ± 0.0108 & 0.4760 ± 0.2277 \\
6  & 0.2312 & 0.4198 & 0.0037 ± 0.0101 & 0.2275 ± 0.0101 & 0.4161 ± 0.0101 & 0.4760 ± 0.2277 \\
7  & 0.6919 & 0.6896 & 0.0037 ± 0.0139 & 0.6882 ± 0.0139 & 0.6859 ± 0.0139 & 0.4760 ± 0.2277 \\
8  & 0.6824 & 0.5694 & 0.0037 ± 0.0143 & 0.6787 ± 0.0143 & 0.5657 ± 0.0143 & 0.4760 ± 0.2277 \\
9  & 0.6913 & \textbf{0.0903} & 0.0036 ± 0.0061 & 0.6877 ± 0.0061 & \textbf{0.0867 ± 0.0061} & 0.4760 ± 0.2277 \\
10 & \textbf{0.0733} & 0.1276 & 0.0036 ± 0.0206 & \textbf{0.0697 ± 0.0206} & 0.1240 ± 0.0206 & 0.4760 ± 0.2277 \\
11 & 0.5529 & 0.6932 & 0.0035 ± 0.0067 & 0.5494 ± 0.0067 & 0.6897 ± 0.0067 & 0.4761 ± 0.2277 \\
12 & 0.5981 & 0.6909 & 0.0034 ± 0.0098 & 0.5947 ± 0.0098 & 0.6875 ± 0.0098 & 0.4762 ± 0.2277 \\
13 & 0.5111 & 0.6873 & 0.0031 ± 0.0079 & 0.5080 ± 0.0079 & 0.6842 ± 0.0079 & 0.4764 ± 0.2276 \\
14 & 0.4116 & 0.5636 & 0.0025 ± 0.0063 & 0.4091 ± 0.0063 & 0.5611 ± 0.0063 & 0.4767 ± 0.2275 \\
15 & 0.6620 & 0.5387 & 0.0013 ± 0.0054 & 0.6607 ± 0.0054 & 0.5374 ± 0.0054 & 0.4772 ± 0.2271 \\
16 & 0.2525 & 0.6352 & -0.0010 ± 0.0099 & 0.2535 ± 0.0099 & 0.6362 ± 0.0099 & 0.4779 ± 0.2263 \\
17 & 0.3450 & 0.5577 & -0.0059 ± 0.0138 & 0.3509 ± 0.0138 & 0.5636 ± 0.0138 & 0.4788 ± 0.2246 \\
18 & 0.6929 & 0.6884 & -0.0156 ± 0.0142 & 0.7085 ± 0.0142 & 0.7040 ± 0.0142 & 0.4801 ± 0.2209 \\
19 & 0.2151 & 0.2619 & -0.0355 ± 0.0050 & 0.2506 ± 0.0050 & 0.2974 ± 0.0050 & 0.4819 ± 0.2128 \\
\midrule
\multicolumn{1}{r}{\textbf{Minima}} &
\textbf{0.0733} &
\textbf{0.0903} &
-- &
$\hat{\mathcal{R}}^\infty(Q^*) = 0.0697 \pm 0.0206$ &
$\hat{\mathcal{R}}^\infty(Q^*) = 0.0867 \pm 0.0061$ &
-- \\
\bottomrule
\end{tabular}
}
\caption{Estimated risks and $\Lambda(k)$ values for different history lengths $k$ (\textit{M} = 10,000,000).}
\label{tab:universal_learning_4}
\end{table}





\section{Computational Cost}
\label{sec:cost}
All experiments were conducted on a system equipped with a 14-core CPU, a 20-core integrated GPU, 24~GB of unified memory, and 1~TB of SSD storage. 

\textbf{Estimating $\mathbf{I}_{\text{pred}}$} The computational complexity of the proposed algorithm is primarily driven by the estimation of mutual information between learned representations. This process incurs a cost of $\mathcal{O}(B^2 d)$ per iteration, where $B$ is the batch size and $d$ is the dimensionality of the representations. The $\mathcal{O}(B^2 d)$ complexity reflects the need to compute pairwise interactions between samples within each batch in a high-dimensional space. The total computational cost per iteration is therefore $\mathcal{O}(B^2 d)$, excluding the contribution of the encoder $g$ and decoder $h$, whose complexity depends on their specific architectures.

\textbf{Estimating $\hat{\Lambda}(k)$} Using Corollary~\ref{cor:universal_learning_curve_parametric}, when \( k \) ranges from 1 to \( n \), the computational complexity of estimating the learning curve is \( \mathcal{O}(n B^2 d) \).

\textbf{Estimating $\mathcal{R}^{\infty}(Q^*)$} The total cost corresponds to training the model \( Q \) times over \( n \). However, this cost is generally unknown or difficult to quantify directly. Moreover, it must be added to the computational cost required to estimate the learning curve \( \hat{\Lambda}(k) \).

\section{Algorithms training procedures}
\label{alg:evopred}
\begin{algorithm}[H]
\caption{$\mathbf{I}_{\text{pred}}$: Data is sampled in a sequential manner with temporal alignment}
\begin{algorithmic}[1]
\State \textbf{for each training iteration do}
\State \quad Sample $\{(x_i, y_i)\}_{i=1}^B$ from a long sequence $\{z_t\}_{t=1}^N$ such that:
\[
x_i = [z_{t_i - k}, \ldots, z_{t_i - 1}], \quad y_i = [z_{t_i}, \ldots, z_{t_i + k' - 1}]
\]
\State \quad Compute critic scores $S_{ij} = f_\theta(x_i, y_j)$
\State \quad Compute $\mathbf{I}_{\text{pred}}$ using Eq.~(3):
\[
\mathbf{I}_{\text{pred}, i}(k, k')
 := \frac{1}{B} \sum_{i=1}^B f_\theta(x_i, y_i)
- \ln \left( \frac{1}{B(B-1)} \sum_{\substack{i,j=1 \\ i \ne j}}^B \exp(f_\theta(x_i, y_j)) \right)
\]
\State \quad Update critic parameters $\theta$ by maximizing $\mathbf{I}_{\text{pred}}$
\State \textbf{end for}
\end{algorithmic}
\end{algorithm}

% \section{Limitations}
% \label{sec:limitations}
% Estimating $\mathbf{I}_{\text{pred}}$ remains a challenging task in practice. Since the learning curve is derived directly from this quantity, any noise or bias in the estimation can lead to significant errors in the assessment of the optimal achievable risk. In particular, the tendency of current estimators to systematically underestimate $\mathbf{I}_{\text{pred}}$ can misrepresent the true predictive structure of the data. Future research should focus on improving both the computational efficiency and estimation accuracy of $\mathbf{I}_{\text{pred}}$, and on extending the applicability of the framework to real-world datasets where the complexity and noise levels are higher.
  
% A particularly valuable direction would be to evaluate this theoretical framework on real sequential tasks—such as those in natural language processing—where verifying its practical usefulness would provide strong evidence of its relevance beyond synthetic benchmarks.

\section{Limitations}
\label{sec:limitations}
Estimating $\mathbf{I}_{\text{pred}}$ remains a challenging task in practice. Since the learning curve is derived directly from this quantity, any noise or bias in the estimation can lead to significant errors in the assessment of the optimal achievable risk. In particular, the tendency of current estimators to systematically underestimate $\mathbf{I}_{\text{pred}}$ can misrepresent the true predictive structure of the data. Moreover, estimation can suffer from variance and occasional instability, especially in high-dimensional settings or when the underlying process has weak temporal dependencies. While developing a definitive and robust estimator of $\mathbf{I}_{\text{pred}}$ is not the main objective of this work, we view it as an important and promising direction for future research in its own right.

Another limitation lies in the fact that our estimator of the minimal achievable risk is model-dependent, in contrast to model-free approaches such as \texttt{EvoRate}~\cite{zeng2025towards}. We acknowledge that this dependency introduces potential confounds related to model capacity and training variability. However, our framework is designed with a different purpose: while \texttt{EvoRate} aims to detect the mere presence of temporal structure, our objective is to quantify how much of this structure is captured by a given model class. The core quantity of interest—the learning curve $\Lambda(k)$—remains model-independent and reflects the intrinsic predictability of the data across context lengths. The model-specific risk can then be seen as an added diagnostic, enabling one to assess whether predictive limitations stem from insufficient model capacity or from intrinsic data constraints.

Importantly, in our experiments, the estimated minimal achievable risk was generally stable across different architectures, which suggests that the diagnostic remains robust despite being evaluated within specific model classes. Nonetheless, we emphasize that mutual information estimation in sequential settings is still a young and evolving research area. We hope that our contribution helps lay the theoretical foundation for more robust methods—potentially by shifting focus from direct mutual information estimation to the more stable and interpretable learning curve $\Lambda(k)$.

Finally, an important future direction will be to evaluate the framework on large-scale real-world sequential tasks—such as those in natural language processing—where validating its usefulness beyond synthetic benchmarks would provide strong evidence of its practical relevance.


%%%%%%%%%%%%%%%%%%%%%%%%%%%%%%%%%%%%%%%%%%%%%%%%%%%%%%%%%%%%



\newpage
\section*{NeurIPS Paper Checklist}

%%% BEGIN INSTRUCTIONS %%%
The checklist is designed to encourage best practices for responsible machine learning research, addressing issues of reproducibility, transparency, research ethics, and societal impact. Do not remove the checklist: {\bf The papers not including the checklist will be desk rejected.} The checklist should follow the references and follow the (optional) supplemental material.  The checklist does NOT count towards the page
limit. 

Please read the checklist guidelines carefully for information on how to answer these questions. For each question in the checklist:
\begin{itemize}
    \item You should answer \answerYes{}, \answerNo{}, or \answerNA{}.
    \item \answerNA{} means either that the question is Not Applicable for that particular paper or the relevant information is Not Available.
    \item Please provide a short (1–2 sentence) justification right after your answer (even for NA). 
   % \item {\bf The papers not including the checklist will be desk rejected.}
\end{itemize}

{\bf The checklist answers are an integral part of your paper submission.} They are visible to the reviewers, area chairs, senior area chairs, and ethics reviewers. You will be asked to also include it (after eventual revisions) with the final version of your paper, and its final version will be published with the paper.

The reviewers of your paper will be asked to use the checklist as one of the factors in their evaluation. While "\answerYes{}" is generally preferable to "\answerNo{}", it is perfectly acceptable to answer "\answerNo{}" provided a proper justification is given (e.g., "error bars are not reported because it would be too computationally expensive" or "we were unable to find the license for the dataset we used"). In general, answering "\answerNo{}" or "\answerNA{}" is not grounds for rejection. While the questions are phrased in a binary way, we acknowledge that the true answer is often more nuanced, so please just use your best judgment and write a justification to elaborate. All supporting evidence can appear either in the main paper or the supplemental material, provided in appendix. If you answer \answerYes{} to a question, in the justification please point to the section(s) where related material for the question can be found.

IMPORTANT, please:
\begin{itemize}
    \item {\bf Delete this instruction block, but keep the section heading ``NeurIPS Paper Checklist"},
    \item  {\bf Keep the checklist subsection headings, questions/answers and guidelines below.}
    \item {\bf Do not modify the questions and only use the provided macros for your answers}.
\end{itemize} 
 

%%% END INSTRUCTIONS %%%


\begin{enumerate}

\item {\bf Claims}
    \item[] Question: Do the main claims made in the abstract and introduction accurately reflect the paper's contributions and scope?
    \answerYes{}
    \item[] Justification: The claims are clearly stated in the abstract and introduction.
    \item[] Guidelines:
    \begin{itemize}
        \item The answer NA means that the abstract and introduction do not include the claims made in the paper.
        \item The abstract and/or introduction should clearly state the claims made, including the contributions made in the paper and important assumptions and limitations. A No or NA answer to this question will not be perceived well by the reviewers. 
        \item The claims made should match theoretical and experimental results, and reflect how much the results can be expected to generalize to other settings. 
        \item It is fine to include aspirational goals as motivation as long as it is clear that these goals are not attained by the paper. 
    \end{itemize}

\item {\bf Limitations}
    \item[] Question: Does the paper discuss the limitations of the work performed by the authors?
    \item[] Answer: \answerYes{} % Replace by \answerYes{}, \answerNo{}, or \answerNA{}.
    \item[] Justification : See Appendix section ~\ref{sec:limitations}.
    \item[] Guidelines:
    \begin{itemize}
        \item The answer NA means that the paper has no limitation while the answer No means that the paper has limitations, but those are not discussed in the paper. 
        \item The authors are encouraged to create a separate "Limitations" section in their paper.
        \item The paper should point out any strong assumptions and how robust the results are to violations of these assumptions (e.g., independence assumptions, noiseless settings, model well-specification, asymptotic approximations only holding locally). The authors should reflect on how these assumptions might be violated in practice and what the implications would be.
        \item The authors should reflect on the scope of the claims made, e.g., if the approach was only tested on a few datasets or with a few runs. In general, empirical results often depend on implicit assumptions, which should be articulated.
        \item The authors should reflect on the factors that influence the performance of the approach. For example, a facial recognition algorithm may perform poorly when image resolution is low or images are taken in low lighting. Or a speech-to-text system might not be used reliably to provide closed captions for online lectures because it fails to handle technical jargon.
        \item The authors should discuss the computational efficiency of the proposed algorithms and how they scale with dataset size.
        \item If applicable, the authors should discuss possible limitations of their approach to address problems of privacy and fairness.
        \item While the authors might fear that complete honesty about limitations might be used by reviewers as grounds for rejection, a worse outcome might be that reviewers discover limitations that aren't acknowledged in the paper. The authors should use their best judgment and recognize that individual actions in favor of transparency play an important role in developing norms that preserve the integrity of the community. Reviewers will be specifically instructed to not penalize honesty concerning limitations.
    \end{itemize}

\item {\bf Theory assumptions and proofs}
    \item[] Question: For each theoretical result, does the paper provide the full set of assumptions and a complete (and correct) proof?
    \item[] Answer: \answerYes{} % Replace by \answerYes{}, \answerNo{}, or \answerNA{}.
    \item[] Justification: The full set of assumptions is described before each proposition and theorem and each corresponding proof can be found in the Appendix section ~\ref{appendix:partA}.
    \item[] Guidelines:
    \begin{itemize}
        \item The answer NA means that the paper does not include theoretical results. 
        \item All the theorems, formulas, and proofs in the paper should be numbered and cross-referenced.
        \item All assumptions should be clearly stated or referenced in the statement of any theorems.
        \item The proofs can either appear in the main paper or the supplemental material, but if they appear in the supplemental material, the authors are encouraged to provide a short proof sketch to provide intuition. 
        \item Inversely, any informal proof provided in the core of the paper should be complemented by formal proofs provided in appendix or supplemental material.
        \item Theorems and Lemmas that the proof relies upon should be properly referenced. 
    \end{itemize}

    \item {\bf Experimental result reproducibility}
    \item[] Question: Does the paper fully disclose all the information needed to reproduce the main experimental results of the paper to the extent that it affects the main claims and/or conclusions of the paper (regardless of whether the code and data are provided or not)?
    \item[] Answer: \answerYes{} % Replace by \answerYes{}, \answerNo{}, or \answerNA{}.
    \item[] Justification: See Appendix sections ~\ref{appendix:partB} and ~\ref{sec:experiment}.
    \item[] Guidelines:
    \begin{itemize}
        \item The answer NA means that the paper does not include experiments.
        \item If the paper includes experiments, a No answer to this question will not be perceived well by the reviewers: Making the paper reproducible is important, regardless of whether the code and data are provided or not.
        \item If the contribution is a dataset and/or model, the authors should describe the steps taken to make their results reproducible or verifiable. 
        \item Depending on the contribution, reproducibility can be accomplished in various ways. For example, if the contribution is a novel architecture, describing the architecture fully might suffice, or if the contribution is a specific model and empirical evaluation, it may be necessary to either make it possible for others to replicate the model with the same dataset, or provide access to the model. In general. releasing code and data is often one good way to accomplish this, but reproducibility can also be provided via detailed instructions for how to replicate the results, access to a hosted model (e.g., in the case of a large language model), releasing of a model checkpoint, or other means that are appropriate to the research performed.
        \item While NeurIPS does not require releasing code, the conference does require all submissions to provide some reasonable avenue for reproducibility, which may depend on the nature of the contribution. For example
        \begin{enumerate}
            \item If the contribution is primarily a new algorithm, the paper should make it clear how to reproduce that algorithm.
            \item If the contribution is primarily a new model architecture, the paper should describe the architecture clearly and fully.
            \item If the contribution is a new model (e.g., a large language model), then there should either be a way to access this model for reproducing the results or a way to reproduce the model (e.g., with an open-source dataset or instructions for how to construct the dataset).
            \item We recognize that reproducibility may be tricky in some cases, in which case authors are welcome to describe the particular way they provide for reproducibility. In the case of closed-source models, it may be that access to the model is limited in some way (e.g., to registered users), but it should be possible for other researchers to have some path to reproducing or verifying the results.
        \end{enumerate}
    \end{itemize}


\item {\bf Open access to data and code}
    \item[] Question: Does the paper provide open access to the data and code, with sufficient instructions to faithfully reproduce the main experimental results, as described in supplemental material?
    \item[] Answer: \answerYes{} % Replace by \answerYes{}, \answerNo{}, or \answerNA{}.
    \item[] Justification : The GitHub link is provided in the introduction.
    \item[] Guidelines:
    \begin{itemize}
        \item The answer NA means that paper does not include experiments requiring code.
        \item Please see the NeurIPS code and data submission guidelines (\url{https://nips.cc/public/guides/CodeSubmissionPolicy}) for more details.
        \item While we encourage the release of code and data, we understand that this might not be possible, so “No” is an acceptable answer. Papers cannot be rejected simply for not including code, unless this is central to the contribution (e.g., for a new open-source benchmark).
        \item The instructions should contain the exact command and environment needed to run to reproduce the results. See the NeurIPS code and data submission guidelines (\url{https://nips.cc/public/guides/CodeSubmissionPolicy}) for more details.
        \item The authors should provide instructions on data access and preparation, including how to access the raw data, preprocessed data, intermediate data, and generated data, etc.
        \item The authors should provide scripts to reproduce all experimental results for the new proposed method and baselines. If only a subset of experiments are reproducible, they should state which ones are omitted from the script and why.
        \item At submission time, to preserve anonymity, the authors should release anonymized versions (if applicable).
        \item Providing as much information as possible in supplemental material (appended to the paper) is recommended, but including URLs to data and code is permitted.
    \end{itemize}


\item {\bf Experimental setting/details}
    \item[] Question: Does the paper specify all the training and test details (e.g., data splits, hyperparameters, how they were chosen, type of optimizer, etc.) necessary to understand the results?
    \item[] Answer: \answerYes{} % Replace by \answerYes{}, \answerNo{}, or \answerNA{}.
    \item[] Justification: See Appendix sections ~\ref{appendix:partB} and ~\ref{sec:experiment}. 
    \item[] Guidelines:
    \begin{itemize}
        \item The answer NA means that the paper does not include experiments.
        \item The experimental setting should be presented in the core of the paper to a level of detail that is necessary to appreciate the results and make sense of them.
        \item The full details can be provided either with the code, in appendix, or as supplemental material.
    \end{itemize}

\item {\bf Experiment statistical significance}
    \item[] Question: Does the paper report error bars suitably and correctly defined or other appropriate information about the statistical significance of the experiments?
    \item[] Answer: \answerYes{} % Replace by \answerYes{}, \answerNo{}, or \answerNA{}.
    \item[] Justification: Confidence intervals are reported whenever computationally feasible. See Section~\ref{sec:Experiment} for details and Appendix~\ref{sec:partC} for additional results.

    \item[] Guidelines:
    \begin{itemize}
        \item The answer NA means that the paper does not include experiments.
        \item The authors should answer "Yes" if the results are accompanied by error bars, confidence intervals, or statistical significance tests, at least for the experiments that support the main claims of the paper.
        \item The factors of variability that the error bars are capturing should be clearly stated (for example, train/test split, initialization, random drawing of some parameter, or overall run with given experimental conditions).
        \item The method for calculating the error bars should be explained (closed form formula, call to a library function, bootstrap, etc.)
        \item The assumptions made should be given (e.g., Normally distributed errors).
        \item It should be clear whether the error bar is the standard deviation or the standard error of the mean.
        \item It is OK to report 1-sigma error bars, but one should state it. The authors should preferably report a 2-sigma error bar than state that they have a 96\% CI, if the hypothesis of Normality of errors is not verified.
        \item For asymmetric distributions, the authors should be careful not to show in tables or figures symmetric error bars that would yield results that are out of range (e.g. negative error rates).
        \item If error bars are reported in tables or plots, The authors should explain in the text how they were calculated and reference the corresponding figures or tables in the text.
    \end{itemize}

\item {\bf Experiments compute resources}
    \item[] Question: For each experiment, does the paper provide sufficient information on the computer resources (type of compute workers, memory, time of execution) needed to reproduce the experiments?
    \item[] Answer: \answerYes{} % Replace by \answerYes{}, \answerNo{}, or \answerNA{}.
    \item[] Justification: See Appendix section ~\ref{sec:cost}.
    \item[] Guidelines:
    \begin{itemize}
        \item The answer NA means that the paper does not include experiments.
        \item The paper should indicate the type of compute workers CPU or GPU, internal cluster, or cloud provider, including relevant memory and storage.
        \item The paper should provide the amount of compute required for each of the individual experimental runs as well as estimate the total compute. 
        \item The paper should disclose whether the full research project required more compute than the experiments reported in the paper (e.g., preliminary or failed experiments that didn't make it into the paper). 
    \end{itemize}
    
\item {\bf Code of ethics}
    \item[] Question: Does the research conducted in the paper conform, in every respect, with the NeurIPS Code of Ethics \url{https://neurips.cc/public/EthicsGuidelines}?
    \item[] Answer: \answerYes{} % Replace by \answerYes{}, \answerNo{}, or \answerNA{}.
    \item[] Justification: Anonymity preserved. 
    \item[] Guidelines:
    \begin{itemize}
        \item The answer NA means that the authors have not reviewed the NeurIPS Code of Ethics.
        \item If the authors answer No, they should explain the special circumstances that require a deviation from the Code of Ethics.
        \item The authors should make sure to preserve anonymity (e.g., if there is a special consideration due to laws or regulations in their jurisdiction).
    \end{itemize}


\item {\bf Broader impacts}
    \item[] Question: Does the paper discuss both potential positive societal impacts and negative societal impacts of the work performed?
    Answer: \answerNo{} % Replace by \answerYes{}, \answerNo{}, or \answerNA{}.
    \item[] Justification: This work constitutes foundational research in machine learning rather than an applied study, and as such, it does not pose any foreseeable negative societal impact.
    \item[] Guidelines:
    \begin{itemize}
        \item The answer NA means that there is no societal impact of the work performed.
        \item If the authors answer NA or No, they should explain why their work has no societal impact or why the paper does not address societal impact.
        \item Examples of negative societal impacts include potential malicious or unintended uses (e.g., disinformation, generating fake profiles, surveillance), fairness considerations (e.g., deployment of technologies that could make decisions that unfairly impact specific groups), privacy considerations, and security considerations.
        \item The conference expects that many papers will be foundational research and not tied to particular applications, let alone deployments. However, if there is a direct path to any negative applications, the authors should point it out. For example, it is legitimate to point out that an improvement in the quality of generative models could be used to generate deepfakes for disinformation. On the other hand, it is not needed to point out that a generic algorithm for optimizing neural networks could enable people to train models that generate Deepfakes faster.
        \item The authors should consider possible harms that could arise when the technology is being used as intended and functioning correctly, harms that could arise when the technology is being used as intended but gives incorrect results, and harms following from (intentional or unintentional) misuse of the technology.
        \item If there are negative societal impacts, the authors could also discuss possible mitigation strategies (e.g., gated release of models, providing defenses in addition to attacks, mechanisms for monitoring misuse, mechanisms to monitor how a system learns from feedback over time, improving the efficiency and accessibility of ML).
    \end{itemize}
    
\item {\bf Safeguards}
    \item[] Question: Does the paper describe safeguards that have been put in place for responsible release of data or models that have a high risk for misuse (e.g., pretrained language models, image generators, or scraped datasets)?
    \item[] Answer: \answerNA{} % Replace by \answerYes{}, \answerNo{}, or \answerNA{}.
    \item[] Justification: No such risks are associated with this paper.
    \item[] Guidelines:
    \begin{itemize}
        \item The answer NA means that the paper poses no such risks.
        \item Released models that have a high risk for misuse or dual-use should be released with necessary safeguards to allow for controlled use of the model, for example by requiring that users adhere to usage guidelines or restrictions to access the model or implementing safety filters. 
        \item Datasets that have been scraped from the Internet could pose safety risks. The authors should describe how they avoided releasing unsafe images.
        \item We recognize that providing effective safeguards is challenging, and many papers do not require this, but we encourage authors to take this into account and make a best faith effort.
    \end{itemize}

\item {\bf Licenses for existing assets}
    \item[] Question: Are the creators or original owners of assets (e.g., code, data, models), used in the paper, properly credited and are the license and terms of use explicitly mentioned and properly respected?
    \item[] Answer: \answerYes{} % Replace by \answerYes{}, \answerNo{}, or \answerNA{}.
    \item[] Justification: Relevant references are appropriately acknowledged in the Related Work section.
    \item[] Guidelines:
    \begin{itemize}
        \item The answer NA means that the paper does not use existing assets.
        \item The authors should cite the original paper that produced the code package or dataset.
        \item The authors should state which version of the asset is used and, if possible, include a URL.
        \item The name of the license (e.g., CC-BY 4.0) should be included for each asset.
        \item For scraped data from a particular source (e.g., website), the copyright and terms of service of that source should be provided.
        \item If assets are released, the license, copyright information, and terms of use in the package should be provided. For popular datasets, \url{paperswithcode.com/datasets} has curated licenses for some datasets. Their licensing guide can help determine the license of a dataset.
        \item For existing datasets that are re-packaged, both the original license and the license of the derived asset (if it has changed) should be provided.
        \item If this information is not available online, the authors are encouraged to reach out to the asset's creators.
    \end{itemize}

\item {\bf New assets}
    \item[] Question: Are new assets introduced in the paper well documented and is the documentation provided alongside the assets?
    \item[] Answer: \answerYes{} % Replace by \answerYes{}, \answerNo{}, or \answerNA{}.
    \item[] Justification: Documentation of the new assets can be found in Appendix sections \ref{appendix:partB} and ~\ref{sec:experiment}.
    \item[] Guidelines:
    \begin{itemize}
        \item The answer NA means that the paper does not release new assets.
        \item Researchers should communicate the details of the dataset/code/model as part of their submissions via structured templates. This includes details about training, license, limitations, etc. 
        \item The paper should discuss whether and how consent was obtained from people whose asset is used.
        \item At submission time, remember to anonymize your assets (if applicable). You can either create an anonymized URL or include an anonymized zip file.
    \end{itemize}

\item {\bf Crowdsourcing and research with human subjects}
    \item[] Question: For crowdsourcing experiments and research with human subjects, does the paper include the full text of instructions given to participants and screenshots, if applicable, as well as details about compensation (if any)? 
    \item[] Answer: \answerNA{} % Replace by \answerYes{}, \answerNo{}, or \answerNA{}.
    \item[] Justification: This paper does not involve crowdsourcing or research involving human subjects.
    \item[] Guidelines:
    \begin{itemize}
        \item The answer NA means that the paper does not involve crowdsourcing nor research with human subjects.
        \item Including this information in the supplemental material is fine, but if the main contribution of the paper involves human subjects, then as much detail as possible should be included in the main paper. 
        \item According to the NeurIPS Code of Ethics, workers involved in data collection, curation, or other labor should be paid at least the minimum wage in the country of the data collector. 
    \end{itemize}

\item {\bf Institutional review board (IRB) approvals or equivalent for research with human subjects}
    \item[] Question: Does the paper describe potential risks incurred by study participants, whether such risks were disclosed to the subjects, and whether Institutional Review Board (IRB) approvals (or an equivalent approval/review based on the requirements of your country or institution) were obtained?
    \item[] Answer: \answerNA{} % Replace by \answerYes{}, \answerNo{}, or \answerNA{}.
    \item[] Justification: This paper does not involve crowdsourcing or research involving human subjects.
    \item[] Guidelines:
    \begin{itemize}
        \item The answer NA means that the paper does not involve crowdsourcing nor research with human subjects.
        \item Depending on the country in which research is conducted, IRB approval (or equivalent) may be required for any human subjects research. If you obtained IRB approval, you should clearly state this in the paper. 
        \item We recognize that the procedures for this may vary significantly between institutions and locations, and we expect authors to adhere to the NeurIPS Code of Ethics and the guidelines for their institution. 
        \item For initial submissions, do not include any information that would break anonymity (if applicable), such as the institution conducting the review.
    \end{itemize}

\item {\bf Declaration of LLM usage}
    \item[] Question: Does the paper describe the usage of LLMs if it is an important, original, or non-standard component of the core methods in this research? Note that if the LLM is used only for writing, editing, or formatting purposes and does not impact the core methodology, scientific rigorousness, or originality of the research, declaration is not required.
    %this research? 
    \item[] Answer: \answerNA{} % Replace by \answerYes{}, \answerNo{}, or \answerNA{}.
    \item[] Justification: LLM are only used for editing and formatting purposes.
    \item[] Guidelines:
    \begin{itemize}
        \item The answer NA means that the core method development in this research does not involve LLMs as any important, original, or non-standard components.
        \item Please refer to our LLM policy (\url{https://neurips.cc/Conferences/2025/LLM}) for what should or should not be described.
    \end{itemize}

\end{enumerate}


\end{document}
